\SourcePage{392}
\SGAchapter{XIX}{等特征优良预概形的上同调}{Cohomologie des préschémas excellents d’égales caractéristiques}
\begin{center}M. ARTIN\end{center}
本讲在承认下述\Term{form}{形式}的\Term{resolution-singularities}{奇点消解}定理的前提下，发展\Term{equal-characteristic}{等特征}\Term{excellent}{优良}\Term{prescheme}{预概形}（EGA IV 7.8.1）的\Term{etale}{平展}\Term{cohomology}{上同调}理论：
设 $X=\operatorname{Spec}A$ 是\Term{equal-characteristic}{等特征}的\Term{excellent}{优良}\Term{local-scheme}{局部}\Term{scheme}{概形}，$U\subset X$ 是\Term{regular}{正则}\Term{open-subset}{开集}。存在\Term{proper}{固有}\Term{birational}{双有理}\Term{morphism}{态射} $f:X'\to X$ 和在 $X$ 上的\Term{open-immersion}{开浸入} $U\to X'$，使 $X'$ \Term{regular}{正则}，且 $Y=X-U$\footnote{此处原文及扫描补源均写作 $X-U$，没有撇号。} 处处\Term{codimension}{余维}为 $1$，并具有横截交。
因此，这些证明“目前”只在\Term{characteristic-zero}{特征零}时\Term{effective}{有效} [4]。利用它们也可推出\Term{characteristic}{特征} $p>0$ 的\Term{prescheme}{预概形}在低维时的一些结果 [1]。
在\Term{resolution}{消解}假设下，对于\Term{equal-characteristic}{等特征}\Term{prescheme}{预概形}以及与\Term{characteristic}{特征}互素的系数，我们将得到较为满意的结果（见 3.2、4.1、5.1、6.1）。相反，对于一般混合\Term{characteristic}{特征}\Term{prescheme}{预概形}的\Term{cohomology}{上同调}，我们几乎一无所知；即使在已有\Term{resolution}{消解}的二维情形 [2]，也是如此。例如，
\SourcePage{393}
对于任何\Term{dimension}{维数}大于 $1$ 的混合\Term{characteristic}{特征}\Term{complete}{完备}环，都还不知道\Term{purity}{纯性}定理 2.1。对于\Term{formal-power-series-ring}{形式幂级数环} $k[[x_1,\ldots,x_n]]$ 的这一结果（1.2）是本理论的主要工具之一。另外应指出，1.2 的证明不用\Term{resolution}{消解}定理，因而该定理已在任意\Term{characteristic}{特征}下得到证明。
我们将经常不加说明地使用 EGA IV 7.8 中发展的\Term{excellent}{优良}\Term{scheme}{概形}的性质，如在\Term{finite-type}{有限型}\Term{group-extension}{扩张}下的\Term{stable}{稳定}性（EGA IV 7.8.3(ii)），以及在\Term{strict-localization}{严格局部化}下的\Term{stable}{稳定}性。后者由\Term{resolution}{消解}及 EGA IV 7.9.5 得出。
\SGAsection{1}{环 $k[[x_1,\ldots,x_n]]$ 的纯性}
回顾 XVI 3.7 的下述推论。
\begin{SGAresult}{推论}{1.1}
设 $A$ 是\Term{strictly-local}{严格局部}环，以 $A\{x\}$ 记 $A[x]$ 在由 $\operatorname{rad}A$ 和 $x$ 生成的\Term{prime-ideal}{素理想}处的\Term{strict-localization}{严格局部化}。设 $U=\operatorname{Spec}A\{x\}[1/x]$。则对与 $A$ 的\Term{residue-characteristic}{剩余特征}互素的 $n$，有
\[H^q(U,\mu_n)\simeq\begin{cases}\mu_n(A)&q=0,\\\mathbf Z/n&q=1,\\0&q>1.\end{cases}\]
\end{SGAresult}
事实上，令 $X=\operatorname{Spec}A[x]$，并以 $Y=V(x)$ 记 $x$ 在 $X$ 中的零点集。于是 $(Y,X)$ 是 $\operatorname{Spec}A$ 上的\Term{smooth}{光滑}对（XVI 3.1）。据 V 4、VIII 5，$H^q(U,\mu_n)$ 等同于\Term{sheaf}{层} $\mathcal H_Y^{q+1}(X,\mu_n)$ 在 $X$ 的点 $(\operatorname{rad}A,x)$ 上方某个\Term{geometric-point}{几何点}处的\Term{fiber}{纤维}，故推论直接由 XVI 3.7 得出。注意 $\mu_n$ 和 $\mathbf Z/n$ 都是 $n$ 阶循环群；这里将它们写成上述\Term{form}{形式}，是为了给出\Term{canonical}{典范}公式（见 3.4）。
\begin{SGAresult}{定理}{1.2}
设 $k$ 是\Term{separably-closed-field}{可分闭域}，$U=\operatorname{Spec}k[[x_1,\ldots,x_r]][1/x_r]$。则对与 $k$ 的\Term{characteristic}{特征}互素的 $n$，有
\[H^q(U,\mu_n)=\begin{cases}\mu_n(A)&q=0,\\\mathbf Z/n&q=1,\\0&q>1.\end{cases}\]
\end{SGAresult}
\textbf{注记。}当 $k=\mathbf C$ 时，证明也适用于收敛幂级数环。此外，只要更仔细地表述结论，就容易去掉 $k$ \Term{separably-closed}{可分闭}的假设。但为了在\Term{resolution-singularities}{奇点消解}假设下建立一般\Term{purity}{纯性}定理 3.2，我们只需\Term{separably-closed}{可分闭}情形；其余断言则是 3.2 的推论。
\SourcePage{394}
对于 $0$、$1$ 维\Term{cohomology}{上同调}，可对任意\Term{regular}{正则}\Term{strictly-local}{严格局部}环证明这一结果。
\begin{SGAresult}{引理}{1.3}
设 $A$ 是\Term{regular}{正则}\Term{strictly-local}{严格局部}环，$x\in\operatorname{rad}A$ 是\Term{local-scheme}{局部}参数，$U=\operatorname{Spec}A[1/x]$。则 $H^0(U,\mu_n)=\mu_n(A)$，且有\Term{canonical}{典范}\Term{isomorphism}{同构}
\[\mathbf Z/n\simeq H^0(U,\mathbf G_m)/(H^0(U,\mathbf G_m))^n\simeq H^1(U,\mu_n),\]
其中 $\mathbf Z/n$ 的\Term{generator}{生成元}是 $x\in H^0(U,\mathbf G_m)$ 的“剩余类”。
\end{SGAresult}
\begin{proof}
$q=0$ 时，断言就是 $U$ \Term{nonempty}{非空}且\Term{connected}{连通}，这确实成立。$q=1$ 时，应用 Kummer \Term{exact-sequence}{正合列}（IX 3.7）。因为 $A$ \Term{regular}{正则}且为\Term{factorial}{析因}环，而 $U$ 是 $X=\operatorname{Spec}A$ 的\Term{open-subset}{开集}，所以 $\operatorname{Pic}U=0$。由此有\Term{exact-sequence}{正合列}
\[H^0(U,\mathbf G_m)\xrightarrow{n}H^0(U,\mathbf G_m)\longrightarrow H^1(U,\mu_n)\longrightarrow0.\]
只需证明 $H^0(U,\mathbf G_m)/(H^0(U,\mathbf G_m))^n$ \Term{isomorphism}{同构}于 $\mathbf Z/n$，并由 $x$ 的剩余类生成。显然有\Term{exact-sequence}{正合列}
\[\begin{tikzcd}0\arrow[r]&H^0(X,\mathbf G_m)\arrow[r]\arrow[d,"\simeq"]&H^0(U,\mathbf G_m)\arrow[r]\arrow[d,"\simeq"]&\mathbf Z\arrow[r]&0\\&A^*&A[1/x]^*\end{tikzcd}\]
其中\Term{section}{截面} $a\in H^0(U,\mathbf G_m)$ 在 $\mathbf Z$ 中的像是它沿 $\{x=0\}$ 的零点阶数。群 $A^*$ 是 $n$ 可除的，因为 $A$ \Term{strictly-local}{严格局部}，且 $n$ 与其\Term{residue-characteristic}{剩余特征}互素。事实上，对 $u\in H^0(X,\mathbf G_m)$，方程 $Y^n-u=0$ 定义 $A$ 的\Term{finite}{有限}\Term{etale}{平展}\Term{group-extension}{扩张}，因而\Term{completely-split}{完全分裂}。故 $H^1(U,\mu_n)\simeq\mathbf Z/n$，引理得证。
\end{proof}
\textbf{1.2 的证明。}只剩 $q>1$ 的情形。回顾 V 附录中的术语。有
\[H^q(U,\mu_n)=\varinjlim_{\mathcal U}H^q(\mathcal U,\mu_n),\]
其中 $\mathcal U=(U_0\leftleftarrows U_1\mathrel{\substack{\leftarrow\\[-.6ex]\leftarrow\\[-.6ex]\leftarrow}}\cdots)$ 遍历 $U$ 的\Term{hypercovering}{超覆盖}\Term{category}{范畴}（在\Term{homotopy}{同伦}意义下\Term{filtered}{滤过}）。取这样一个\Term{hypercovering}{超覆盖} $\mathcal U$，并令各 $U_i$ \Term{separated}{分离}且为\Term{finite-type}{有限型}。为了证明 $H^q(U,\mu_n)=0$，只需找到\Term{hypercovering}{超覆盖} $\mathcal V$ 及\Term{morphism}{态射} $\mathcal V\to\mathcal U$，使诱导\Term{morphism}{态射}
\[H^q(\mathcal U,\mu_n)\longrightarrow H^q(\mathcal V,\mu_n)\]
为零。设 $\{W_1,\ldots,W_m\}$ 是所有 $i\le q+1$ 的 $U_i$ 的\Term{connected-component}{连通分支}之集，$B_j$ 是 $A=k[[x_1,\ldots,x_r]]$ 在 $W_j$ 的\Term{rational-function-field}{有理函数域} $R(W_j)$ 中的\Term{normalization}{正规化}。由于 $W_j\to U$ \Term{etale}{平展}，$W_j$
\SourcePage{395}
是 $\operatorname{Spec}B_j$ 的\Term{open-subset}{开集}，设其补集为 $\operatorname{Spec}C_j$，其中 $C_j$ 为 $B_j$ 的适当\Term{reduced}{约化}\Term{quotient-ring}{商环}。
再取 $b_j\in B_j$，使它生成 $k((x_1,\ldots,x_r))$ 的\Term{separable-extension}{可分扩张} $R(W_j)$。设 $f_j(Y)\in A[Y]$ 为 $b_j$ 在 $A$ 上的首一\Term{irreducible}{不可约}多项式，故 $B_j$ \Term{isomorphism}{同构}于 $A[Y]/(f_j)$ 的\Term{normalization}{正规化}。最后，令 $d_j\in A$ 为 $f_j(Y)$ 的\Term{discriminant}{判别式}，它是非零元。
我们可以把 $i<r$ 的 $x_i$ 换为 $A$ 的其他元素，唯一条件是 $\{x_1,\ldots,x_r\}$ 仍构成 $A$ 的\Term{system-of-parameters}{参数系}。容易看出，必要时作这样的改变，可假定满足下列条件。
\textbf{1.4。}(i) 对每个 $d_j$ 不是单位的 $j$，$x_1,\ldots,x_{r-1},d_j$ 生成一个以 $\operatorname{rad}A$ 为根的\Term{primary}{准素}\Term{ideal}{理想}。
(ii) 对每个 $j$，$x_1,\ldots,x_{r-1}$ 在 $C_j$ 中的像生成一个以 $\operatorname{rad}C_j$ 为根的\Term{primary}{准素}\Term{ideal}{理想}。
\begin{SGAresult}{引理}{1.5}
令 $A'=k[[x_1,\ldots,x_{r-1}]]\{x_r\}$ 为 $k[[x_1,\ldots,x_{r-1}]][x_r]$ 在点 $(x_1,\ldots,x_r)$ 处的 Hensel 化，并令 $X'=\operatorname{Spec}A'$，$U'=\operatorname{Spec}A'[1/x_r]$。在上述条件 1.4 下，存在 $U'$ 的一个（$q+1$ 阶）\Term{hypercovering}{超覆盖} $\mathcal U'$ 及\Term{isomorphism}{同构} $\mathcal U'\times_{X'}X\simeq\mathcal U$。
\end{SGAresult}
暂承认引理。环 $A'$ \Term{excellent}{优良}（EGA IV 7.8.3(ii)、7.9.5），故 $U'_i$ 和 $U_i$ 的\Term{connected-component}{连通分支}一一对应（EGA IV 7.8.3(vii)）。因此上述\Term{isomorphism}{同构}给出
\[H^q(\mathcal U',\mu_n)\simeq H^q(\mathcal U,\mu_n).\]
因 $k$ \Term{separably-closed}{可分闭}，$A'$ \Term{strictly-local}{严格局部}，可对 $A'$ 应用 1.1，得到 $H^q(U',\mu_n)=0$。但
\[H^q(U',\mu_n)\simeq\varinjlim_{\mathcal V'}H^q(\mathcal V',\mu_n),\]
其中 $\mathcal V'$ 遍历 $U'$ 的 $q+1$ 阶\Term{hypercovering}{超覆盖}\Term{category}{范畴}，且 $H^q(\mathcal U',\mu_n)$ 是\Term{finite-group}{有限群}。因此存在\Term{hypercovering}{超覆盖}\Term{morphism}{态射} $\mathcal V'\to\mathcal U'$，使诱导\Term{morphism}{态射}
\[H^q(\mathcal U',\mu_n)\longrightarrow H^q(\mathcal V',\mu_n)\]
为零。令 $\mathcal V=\mathcal V'\times_{X'}X$。则 $H^q(\mathcal V',\mu_n)\simeq H^q(\mathcal V,\mu_n)$（理由与 $\mathcal U'$ 相同），故
\[H^q(\mathcal U,\mu_n)\longrightarrow H^q(\mathcal V,\mu_n)\]
\SourcePage{396}
也为零，证毕。
\begin{proof}[1.5 的证明]
只需找\Term{finite}{有限} $A'$ \Term{algebra}{代数} $B'_j,C'_j$ 以及 $A$ \Term{isomorphism}{同构}
\[B'_j\otimes_{A'}A\xrightarrow{\sim}B_j,\qquad C'_j\otimes_{A'}A\xrightarrow{\sim}C_j.\]
事实上，依 [3] 1.4，\Term{morphism}{态射} $B_j\to C_j$ 由一个\Term{morphism}{态射} $B'_j\to C'_j$ 诱导，故\Term{open-subset}{开集} $W_j\subset\operatorname{Spec}B_j$ 由\Term{open-subset}{开集} $W'_j\subset\operatorname{Spec}B'_j$ 诱导。此外，$W_j$ 间的 $U$ \Term{morphism}{态射}由 $A$ \Term{algebra}{代数} $B_j$ 间的\Term{morphism}{态射}诱导，依 [3] 1.4 又由 $B'_j$ 间的\Term{morphism}{态射}诱导，因而由 $W'_j$ 间的 $U'$ \Term{morphism}{态射}诱导。由此，各 $W'_j$ 构成 $U'$ 上一个\Term{simplicial-object}{单纯对象} $\mathcal U'$ 的\Term{connected-component}{连通分支}之集；此\Term{object}{对象}诱导 $\mathcal U$，且显然为 $U'$ 的\Term{hypercovering}{超覆盖}。
条件 1.4(ii) 确实蕴含 $C_j$ 由一个 $A'$ \Term{algebra}{代数} $C'_j$ 诱导，这是容易验证的。\footnote{事实上，1.4(ii) 蕴含 $C_j$ 在 $k[[x_1,\ldots,x_{r-1}]]$ 上\Term{finite}{有限}，从而在 $A'$ 上\Term{finite}{有限}；只须取 $C'_j=C_j$，其 $A'$ 模结构由\Term{restriction-of-scalars}{限制标量}得到。}对于 $B_j$，只需找\Term{monic-polynomial}{首一多项式} $f'_j(Y)\in A'[Y]$，使两个 $A$ \Term{algebra}{代数}
\[A[Y]/(f_j),\qquad A[Y]/(f'_j)=A\otimes_{A'}A'[Y]/(f'_j)\]
\Term{isomorphism}{同构}。于是 $A'[Y]/(f'_j)$ \Term{reduced}{约化}，取其\Term{normalization}{正规化}为 $B'_j$（EGA IV 7.8.3(vi)）。环 $B'_j\otimes_{A'}A$ 仍\Term{normal}{正规}（同上 (vii)），且与 $B_j$ \Term{birational}{双有理}，故等于 $B_j$。
由 1.4(i)，$d_j$ 生成的\Term{ideal}{理想} $I_j\subset A$ 由 $k[[x_1,\ldots,x_{r-1}]][x_r]$ 的一个\Term{ideal}{理想}诱导（Weierstrass 预备定理 [5][6]），因而由 $A'$ 的\Term{ideal}{理想} $I'_j$ 诱导。此外，$A/I_j$ 是\Term{finite}{有限} $k[[x_1,\ldots,x_{r-1}]]$ \Term{algebra}{代数}，所以\Term{isomorphism}{同构}于 $A'/I'_j$。后者因此\Term{complete}{完备}，故 $A'$ 的 $I'_j$ 进\Term{completion}{完备化}\Term{complete}{完备}且\Term{isomorphism}{同构}于 $A$。于是可取\Term{monic-polynomial}{首一多项式} $f'_j\in A'[Y]$，使在 $A[Y]$ 中
\[f_j\equiv f'_j\pmod{d_j^2\operatorname{rad}A}.\]
由 [3] 1.3，$A[Y]/(f_j)\simeq A[Y]/(f'_j)$，引理得证。
\end{proof}
\SGAsection{2}{严格局部环的情形}
定理如下。
\begin{SGAresult}{定理}{2.1}
\footnote{依赖\Term{resolution-singularities}{奇点消解}（见引言）。}设 $A$ 是\Term{excellent}{优良}、\Term{regular}{正则}、\Term{strictly-local}{严格局部}的\Term{equal-characteristic}{等特征}环，
\SourcePage{397}
$x$ 为 $A$ 的一个“\Term{local-scheme}{局部}参数”\footnote{这里指属于某个\Term{regular}{正则}\Term{system-of-parameters}{参数系}的元素。}。令 $U=\operatorname{Spec}A[1/x]$。则对于与 $A$ 的\Term{characteristic}{特征}互素的 $n$，
\[H^q(U,\mu_n)=\begin{cases}\mu_n(A)&q=0,\\\mathbf Z/n&q=1,\\0&q>1.\end{cases}\]
\end{SGAresult}
证明对 $q$ 归纳。以 $\mathrm P(N)$ 记 2.1 对 $q\le N$ 的断言。由 1.3，$\mathrm P(1)$ 成立。
\begin{SGAresult}{引理}{2.2}
设 $\mathrm P(N)$ 成立。设 $A\to A'$ 是\Term{excellent}{优良}、\Term{regular}{正则}、\Term{strictly-local}{严格局部}\Term{equal-characteristic}{等特征}环的\Term{local-homomorphism}{局部同态}。设 $\{x_1,\ldots,x_r\}$ 是 $A$ 的某个\Term{regular}{正则}\Term{system-of-parameters}{参数系}的一部分，且其在 $A$ 中的像\footnote{此处原文及扫描补源均为 $A$，没有撇号。}仍是某个\Term{regular}{正则}\Term{system-of-parameters}{参数系}的一部分。令
\[U=\operatorname{Spec}A[1/x_1,\ldots,1/x_r],\qquad U'=\operatorname{Spec}A'[1/x'_1,\ldots,1/x'_r].\]
则\Term{canonical}{典范}\Term{morphism}{态射} $H^q(U,\mu_n)\to H^q(U',\mu_n)$ 对每个 $q\le N$ 都是\Term{bijection}{双射}。
\end{SGAresult}
\begin{proof}
\footnote{有一个较短的证明，但要用杯积\Term{form}{形式}：在 $\mathrm P(N)$ 下，证明 $H^*(U,\mathbf Z/n\mathbf Z)$ 在次数 $\le N$ 时是 $H^1(U,\mathbf Z/n\mathbf Z)\simeq(\mathbf Z/n\mathbf Z)^r$ 在 $\mathbf Z/n\mathbf Z$ 上的外\Term{algebra}{代数}，这便把 2.2 归结为熟知的 $q=1$ 情形。这里在 $\mathrm P(N)$ 下所说的计算，例如在 SGA 7 中有明确说明。}
对 $r$ 归纳。$r=0$ 时显然。设 $r>0$，且已对 $r-1$ 证明。令 $B=A/(x)$\footnote{原文在此写作 $x$；随后用 $x_r$ 定义闭补集，记号照录。}，$V=\operatorname{Spec}A[1/x_1,\ldots,1/x_{r-1}]$，$Y=\operatorname{Spec}B[1/x_1,\ldots,1/x_{r-1}]$（仍用 $x_i$ 记它在 $B$ 中的像）。有\Term{open-immersion}{开浸入} $i:U\to V$，其闭补集为由 $x_r$ 定义的 $Y$。
对 $V$ 在 $Y$ 的\Term{geometric-point}{几何点}处的\Term{strictly-local}{严格局部}环应用 $\mathrm P(N)$，得到
\[R^qi_*(\mu_n)_U=\begin{cases}(\mu_n)_V&q=0,\\(\mathbf Z/n)_Y&q=1,\\0&1<q\le N.\end{cases}\tag{2.3}\]
因为这些\Term{strictly-local}{严格局部}环是\Term{excellent}{优良}、\Term{regular}{正则}、\Term{equal-characteristic}{等特征}环，而 $x_r$ 在 $Y$ 的每一点都是\Term{local-scheme}{局部}参数。
\SourcePage{398}
用撇号记由 $A'$ 得出的类似\Term{object}{对象}。得到\Term{spectral-sequence}{谱序列}\Term{morphism}{态射}
\[\begin{tikzcd}E_2^{pq}=H^p(V,R^qi_*\mu_n)\arrow[r,Rightarrow]\arrow[d,"\varphi_2^{pq}"]&H^{p+q}(U,\mu_n)\arrow[d,"\varphi"]\\E_2'{}^{pq}=H^p(V',R^qi'_*\mu_n)\arrow[r,Rightarrow]&H^{p+q}(U',\mu_n).\end{tikzcd}\]
由 1.3 及 (2.3)，\Term{base-change}{基变换} $g:V'\to V$ 显然在 $q\le N$ 时诱导\Term{isomorphism}{同构}
\[g^*R^qi_*\mu_n\xrightarrow{\sim}R^qi'_*\mu_n\]
（只有 $q=0,1$ 时有一点需要验证）。故对\Term{morphism}{态射} $H^q(V,\mu_n)\to H^q(V',\mu_n)$ 应用归纳假设，知 $\varphi_2^{pq}$ 在 $q=0,p\le N$ 时是\Term{isomorphism}{同构}。
同样，可对\Term{ring-homomorphism}{环同态} $B\to B'$ 应用归纳假设：两者均为\Term{regular}{正则}、\Term{excellent}{优良}、\Term{strictly-local}{严格局部}\Term{equal-characteristic}{等特征}环，且 $x_1,\ldots,x_{r-1}$ 是\Term{system-of-parameters}{参数系}的一部分。因此
\[H^q(Y,\mu_n)\xrightarrow{\sim}H^q(Y',\mu_n)\quad(q\le N).\]
因 $(\mu_n)_Y\simeq(\mathbf Z/n)_Y$，$\varphi_2^{pq}$ 在 $p\le N,q=1$ 时也\Term{bijection}{双射}。由于 $1<q\le N$ 时 $R^qi_*\mu_n=R^qi'_*\mu_n=0$，有 $E_2^{pq}=E_2'{}^{pq}=0$。因此在 $1\le m\le N$ 时，\Term{spectral-sequence}{谱序列}\Term{morphism}{态射}化为\Term{exact-sequence}{正合列}\Term{morphism}{态射}
\[\begin{tikzcd}[column sep=small]
H^{m-2}(Y,\mathbf Z/n)\arrow[r]\arrow[d,"a"]&H^m(V,\mu_n)\arrow[r]\arrow[d,"b"]&H^m(U,\mu_n)\arrow[r,"e"]\arrow[d,"c"]&H^{m-1}(Y,\mathbf Z/n)\arrow[r]\arrow[d,"d"]&H^{m+1}(V,\mu_n)\arrow[d]\\
H^{m-2}(Y',\mathbf Z/n)\arrow[r]&H^m(V',\mu_n)\arrow[r]&H^m(U',\mu_n)\arrow[r]&H^{m-1}(Y',\mathbf Z/n)\arrow[r]&H^{m+1}(V',\mu_n).
\end{tikzcd}\]
已知 $a,b,d$ \Term{bijection}{双射}，故 $c$ \Term{injectivity}{单射}。要证明 $c$ \Term{surjectivity}{满射}，只需证明
\[e:H^m(U,\mu_n)\longrightarrow H^{m-1}(Y,\mathbf Z/n)\]
在 $m\le N$ 时\Term{surjectivity}{满射}。
取任意\Term{separably-closed}{可分闭}子域 $k\subset A$，以 $A^\circ=k\{x_1,\ldots,x_r\}$ 记 $k[x_1,\ldots,x_r]$ 在原点的 Hensel 化。则\Term{morphism}{态射} $A^\circ\to A$ 和元素 $x_i$ 满足 2.1 的假设。用显然的记号，有\Term{commutative-square}{交换方块}
\SourcePage{399}
\[\begin{tikzcd}H^m(U^\circ,\mu_n)\arrow[r,"e^\circ"]\arrow[d,"c^\circ"]&H^{m-1}(Y^\circ,\mathbf Z/n)\arrow[d,"d^\circ"]\\H^m(U,\mu_n)\arrow[r,"e"]&H^{m-1}(Y,\mathbf Z/n).
\end{tikzcd}\]
已知 $d^\circ$ \Term{bijection}{双射}，故只须证明 $e^\circ$ \Term{surjectivity}{满射}。于是，略去圆圈上标，我们归结为在 $A=k\{x_1,\ldots,x_r\}$ 时验证\Term{surjectivity}{满射}性。
对此环——它是 $k$ 上\Term{smooth}{光滑}\Term{scheme}{概形}的\Term{limit}{极限}——已在 XVI 3.7 证明\Term{purity}{纯性}定理。因此每个 $q>1$ 都有 $R^qi_*\mu_n=0$，从而对每个 $m\ge1$ 有\Term{exact-sequence}{正合列}
\[H^m(U,\mu_n)\xrightarrow e H^{m-1}(Y,\mathbf Z/n)\longrightarrow H^{m+1}(V,\mu_n)\xrightarrow f H^{m+1}(U,\mu_n).\]
故 $e$ \Term{surjectivity}{满射}\Term{equivalence}{等价}于 $f$ \Term{injectivity}{单射}。
令 $Z=\operatorname{Spec}k\{x_1,\ldots,x_{r-1}\}[1/x_1,\ldots,1/x_{r-1}]$。由环包含给出的\Term{morphism}{态射}
$\operatorname{Spec}k\{x_1,\ldots,x_r\}\to\operatorname{Spec}k\{x_1,\ldots,x_{r-1}\}$ \Term{acyclic}{无上同调}；这就是\Term{smooth}{光滑}\Term{morphism}{态射}的\Term{locally-acyclic}{局部无上同调}性（XVI 1.1）。因 $V$ 是 $Z$ 的\Term{inverse-image}{逆像}，\Term{morphism}{态射}
\[H^q(Z,\mu_n)\longrightarrow H^q(V,\mu_n)\]
对每个 $q$ \Term{bijection}{双射}（XV 1.3）。故只须证明 $H^q(Z,\mu_n)\to H^q(U,\mu_n)$ 对每个 $q$ \Term{injectivity}{单射}。
环 $k\{x_1,\ldots,x_r\}$ 是一族环 $A_\alpha$ 的\Term{inductive-limit}{归纳极限}；这些环在 $k\{x_1,\ldots,x_{r-1}\}[x_r]$ 上\Term{etale}{平展}，且点 $x_1=\cdots=x_r=0$ \Term{lifting}{提升}到这些环。因此 $U$ 是\Term{open-subset}{开集} $U_\alpha=\operatorname{Spec}A_\alpha[1/x_1,\ldots,1/x_r]$ 的\Term{inverse-limit}{投射极限}，且
\[H^q(U,\mu_n)=\varinjlim_\alpha H^q(U_\alpha,\mu_n).\]
由于 $A_\alpha[1/x_r]$ 在 Hensel 环 $k\{x_1,\ldots,x_{r-1}\}$ 上\Term{smooth}{光滑}，且其\Term{closed-fiber}{闭纤维}\Term{nonempty}{非空}，容易看出这个\Term{algebra}{代数}有\Term{section}{截面}。因此 $U_\alpha$ 在 $Z$ 上有\Term{section}{截面}，故 $H^q(Z,\mu_n)\to H^q(U_\alpha,\mu_n)$ 对每个 $q$ \Term{injectivity}{单射}，于是 $H^q(Z,\mu_n)\to H^q(U,\mu_n)$ 也\Term{injectivity}{单射}，引理得证。
\end{proof}
\SourcePage{400}
\begin{SGAresult}{引理}{2.4}
设 $\mathrm P(N)$ 成立。设 $A\to A'$ 是\Term{excellent}{优良}、\Term{strictly-local}{严格局部}\Term{equal-characteristic}{等特征}环（不必\Term{regular}{正则}）的\Term{regular}{正则}\Term{homomorphism}{同态}（EGA IV 6.8.1），$U\subset\operatorname{Spec}A$ 是\Term{regular}{正则}\Term{open-subset}{开集}，$U'=U\times_{\operatorname{Spec}A}\operatorname{Spec}A'$。则对与 $A$ 的\Term{residue-characteristic}{剩余特征}互素的 $n$，\Term{morphism}{态射}
\[H^q(U,\mu_n)\longrightarrow H^q(U',\mu_n)\]
在 $q\le N$ 时\Term{bijection}{双射}。
\end{SGAresult}
\begin{proof}
取 $\operatorname{Spec}A$ 的\Term{resolution-singularities}{奇点消解} $f:X\to\operatorname{Spec}A$，使 $U$ \Term{lifting}{提升}为\Term{open-subset}{开集} $i:U\to X$，且闭集 $Y=X-U$ 纯\Term{codimension}{余维} $1$，具有横截交。用撇号记 $A\to A'$ 诱导的\Term{base-change}{基变换}。\Term{morphism}{态射} $g:X'\to X$ \Term{regular}{正则}（EGA IV 6.8.3），故 $f':X'\to\operatorname{Spec}A'$ 是 $\operatorname{Spec}A$ 的\Term{resolution}{消解}\footnote{此处原文在 $A$ 上没有撇号，照录。}，$i':U'\to X'$ 为\Term{open-immersion}{开浸入}，而 $Y'=X'-U'$ 纯\Term{codimension}{余维} $1$ 且具有横截交。所有这些\Term{scheme}{概形}都\Term{excellent}{优良}。
取 $p'\in X'$，其在 $X$ 中的像为 $p$。由于 $Y$ 在 $X$ 中\Term{transverse-intersection}{横截相交}，存在 $x_1,\ldots,x_r\in\mathcal O_{X,p}$，它们是该环\Term{system-of-parameters}{参数系}的一部分，且 $\operatorname{Spec}\mathcal O_{X,p}\cap Y=V(x_1,\ldots,x_r)$。由 $g$ \Term{regular}{正则}立即得出，$x_i$ 在 $\mathcal O_{X',p'}$ 中的像 $x'_i$ 仍是该环\Term{system-of-parameters}{参数系}的一部分，且 $\operatorname{Spec}\mathcal O_{X',p'}\cap Y'=V(x'_1,\ldots,x'_r)$。
令 $B'$ 为 $\mathcal O_{X',p'}$ 在 $p'$ 上方一个\Term{geometric-point}{几何点} $\bar p$ 处的\Term{strict-localization}{严格局部化}，$B$ 为 $\mathcal O_{X,p}$ 在对应\Term{geometric-point}{几何点}处的\Term{strict-localization}{严格局部化}；这些环\Term{excellent}{优良}（EGA IV 7.8.3(ii)、7.9.5）。令
\[U_{\bar p}=\operatorname{Spec}B[1/x_1,\ldots,1/x_r],\quad U'_{\bar p}=\operatorname{Spec}B'[1/x'_1,\ldots,1/x'_r].\]
由引理 2.2，\Term{morphism}{态射}
\[H^q(U_{\bar p},\mu_n)\longrightarrow H^q(U'_{\bar p},\mu_n)\tag{*}\]
在 $q\le N$ 时\Term{bijection}{双射}。计算\Term{open-immersion}{开浸入} $i,i'$ 的\Term{sheaf}{层} $R^qi_*\mu_n,R^qi'_*\mu_n$：据 VIII 5，它们在相应\Term{geometric-point}{几何点}的\Term{fiber}{纤维}恰为 $H^q(U_{\bar p},\mu_n)$ 和 $H^q(U'_{\bar p},\mu_n)$。故 (*) 蕴含\Term{base-change}{基变换}\Term{morphism}{态射}
\[g^*R^qi_*\mu_n\longrightarrow R^qi'_*\mu_n\tag{**}\]
在 $q\le N$ 时\Term{bijection}{双射}。
由\Term{proper}{固有}\Term{base-change}{基变换}定理（XII 5.5），因 $X,X'$ 在\Term{strictly-local}{严格局部}环谱上\Term{proper}{固有}且\Term{closed-fiber}{闭纤维}相同，对任意\Term{torsion-sheaf}{挠层} $F$ 及每个 $q$ 都有
\[H^q(X,F)\xrightarrow{\sim}H^q(X',g^*F).\]
\SourcePage{401}
因此 (**) 蕴含在\Term{spectral-sequence}{谱序列}\Term{morphism}{态射}
\[\begin{tikzcd}E_2^{pq}=H^p(X,R^qi_*\mu_n)\arrow[d,"\varphi_2^{pq}"]\arrow[r,Rightarrow]&H^{p+q}(U,\mu_n)\arrow[d,"\varphi"]\\E_2'{}^{pq}=H^p(X',R^qi'_*\mu_n)\arrow[r,Rightarrow]&H^{p+q}(U',\mu_n)\end{tikzcd}\]
中，$\varphi_2^{pq}$ 对每个 $q\le N$ \Term{bijection}{双射}。容易看出，这蕴含 $\varphi_\infty^{pq}$ 在 $q\le N,p+q\le N+1$ 时\Term{bijection}{双射}。故\Term{abutment}{终项}\Term{morphism}{态射} $H^m(U,\mu_n)\to H^m(U',\mu_n)$ 在 $m\le N$ 时\Term{bijection}{双射}，引理得证。
\end{proof}
下述引理完成 2.1 的证明。
\begin{SGAresult}{引理}{2.5}
$\mathrm P(N)\Rightarrow\mathrm P(N+1)$。
\end{SGAresult}
\begin{proof}
设 $A$ \Term{strictly-local}{严格局部}、\Term{regular}{正则}、\Term{excellent}{优良}、\Term{equal-characteristic}{等特征}，$x\in\operatorname{rad}A$ 为\Term{local-scheme}{局部}参数，$U=\operatorname{Spec}A[1/x]$。以 $\widehat A$ 记 $A$ 的\Term{completion}{完备化}，$\widehat x$ 记 $x$ 的像，$\widehat U=\operatorname{Spec}\widehat A[1/\widehat x]$，$g:\widehat U\to U$ 为\Term{canonical}{典范}\Term{morphism}{态射}。只须证明
\[g_*(\mu_n)_{\widehat U}\simeq(\mu_n)_U,\qquad R^qg_*\mu_n=0\quad(1\le q\le N).\tag{2.6}\]
事实上，Leray \Term{spectral-sequence}{谱序列}
\[E_2^{pq}=H^p(U,R^qg_*\mu_n)\Longrightarrow H^{p+q}(U',\mu_n)\]
将给出 $q\le N$ 时的\Term{isomorphism}{同构} $H^q(U,\mu_n)\xrightarrow{\sim}H^q(U',\mu_n)$ 以及\Term{exact-sequence}{正合列}
\[0\longrightarrow H^{N+1}(U,\mu_n)\longrightarrow H^{N+1}(U',\mu_n)\longrightarrow H^0(U,R^{N+1}g_*\mu_n).\]
已知 $\mathrm P(1)$ 成立（1.3），所以只须证明 $1<q\le N+1$ 时 $H^q(U,\mu_n)=0$。于是可用 $\widehat A$ 代替 $A$；它\Term{isomorphism}{同构}于\Term{formal-power-series-ring}{形式幂级数环} $k[[x_1,\ldots,x_r]]$，其中 $k$ \Term{separably-closed}{可分闭}且 $\widehat x=x_r$。此情形已在 1.2 处理。
现在验证 (2.6)。$g_*(\mu_n)_{\widehat U}\simeq(\mu_n)_U$ 就是 (2.3) 在 $q=0$ 时的断言。回顾 $R^qg_*\mu_n$ 是\Term{presheaf}{预层} $\mathscr R^q$ 的\Term{adjunction}{伴随}\Term{sheaf}{层}，其中对\Term{etale}{平展} $V\to U$，
\[\mathscr R^q(V)=H^q(\widehat V,\mu_n).\]
只需取\Term{separated}{分离}且\Term{finite-type}{有限型}的 $V$。此时 $V$ 是整\Term{scheme}{概形}之和，故只需考虑整的 $V$。令 $B$ 为 $A$ 在 $\Gamma(V,\mathcal O_V)$ 的\Term{field-of-fractions}{分式域}中的\Term{normalization}{正规化}。则 $B$ 是\Term{finite}{有限} $A$ \Term{algebra}{代数}（EGA IV 7.8.3(vi)），
\SourcePage{402}
而 $V$ 显然是 $\operatorname{Spec}B$ 的\Term{regular}{正则}\Term{open-subset}{开集}，$\widehat V$ 是 $\operatorname{Spec}\widehat B$ 的对应\Term{open-subset}{开集}，其中 $\widehat B=B\otimes_A\widehat A$。于是\Term{open-subset}{开集} $V\subset\operatorname{Spec}B$ 和\Term{morphism}{态射} $B\to\widehat B$ 满足 2.4 的假设，故 $q\le N$ 时
\[\mathscr R^q(V)\simeq H^q(V,\mu_n)\simeq H^q(\widehat V,\mu_n).\]
而\Term{presheaf}{预层} $V\mapsto H^q(V,\mu_n)$ 的\Term{adjunction}{伴随}\Term{sheaf}{层}显然为零，引理得证。
\end{proof}
\SGAsection{3}{纯性}
采用类似 XVI 3 的术语。
\begin{SGAresult}{定义}{3.1}
所谓\Term{codimension}{余维} $c$ 的\Term{regular-pair}{正则对} $(Y,X)$，是\Term{locally-noetherian}{局部 Noether} \Term{regular}{正则}\Term{scheme}{概形}的\Term{closed-immersion}{闭浸入} $i:Y\to X$，使每个 $y\in Y$ 都满足 $\operatorname{codim}_y(Y,X)=c$。此时存在 $y$ 在 $X$ 中的\Term{open-neighborhood}{开邻域} $X'$ 以及\Term{section}{截面} $x_1,\ldots,x_c\in\Gamma(X',\mathcal O_X)$，生成 $Y$ 的\Term{ideal}{理想}，并在 $Y\cap X'$ 每一点构成 $X$ 的\Term{regular}{正则}\Term{system-of-parameters}{参数系}的一部分。记 $U=X-Y$，$j:U\to X$ 为相应\Term{open-immersion}{开浸入}。
\end{SGAresult}
\Term{codimension}{余维} $c$ 的\Term{regular-pair}{正则对}\Term{morphism}{态射} $(Y',X')\to(Y,X)$ 是满足 $Y'=Y\times_XX'$ 的\Term{morphism}{态射} $X'\to X$。\Term{regular-triple}{正则三元组}及其\Term{morphism}{态射}等概念按显然类似的方式定义。
\begin{SGAresult}{定理}{3.2}
（\Term{cohomology}{上同调}\Term{purity}{纯性}）\footnote{依赖\Term{resolution-singularities}{奇点消解}（见引言）。}设 $(Y,X)$ 为\Term{codimension}{余维} $c>0$ 的\Term{regular-pair}{正则对}，其中 $X$ 是\Term{excellent}{优良}\Term{equal-characteristic}{等特征}\Term{prescheme}{预概形}。设 $F$ 是 $X$ 上的\Term{locally-constant}{局部常值}\Term{sheaf}{层}，其群为 $n$ 阶循环群，$n$ 与 $X$ 的\Term{characteristic}{特征}互素。则按 V 4、VIII 5 的记号，
\[\mathcal H_Y^q(X,F)=(R^qi^!)F=0\quad(q\ne2c),\]
\[(R^qj_*)j^*F=0\quad(q\ne0,2c-1),\]
且 $\mathcal H_Y^{2c}(X,F)=i^*(R^{2c-1}j_*)j^*F$ 是 $Y$ 上的\Term{locally-constant}{局部常值} $n$ 阶循环群\Term{sheaf}{层}。
\end{SGAresult}
\begin{proof}[证明（参见 XVI 3.7）]
显然 $q=0,1$ 时 $\mathcal H_Y^q(X,F)=0$。由 V 4.5，
\[\mathcal H_Y^q(X,F)=(R^qi^!)F\simeq i^*(R^{q-1}j_*)j^*F\quad(q>1),\]
故关于这三个\Term{sheaf}{层}的断言\Term{equivalence}{等价}。由于断言对 $X$ 的\Term{etale}{平展}\Term{topology}{拓扑}是\Term{local-scheme}{局部}的，可设 $F=\mu_n$。
\SourcePage{403}
$c=1$ 时，计算 $(R^qj_*)j^*\mu_n$ 的\Term{fiber}{纤维}（VIII 5），归结为 2.1 的情形，因此对 $Y$ 的每个\Term{geometric-point}{几何点} $\bar y$，
\[((R^qj_*)j^*\mu_n)_{\bar y}=\begin{cases}\mathbf Z/n&q=1,\\0&q>1.\end{cases}\]
这些\Term{sheaf}{层}在 $Y$ 外本来为零，故这证明了 $q>1$ 时的断言。$q=1=2c-1$ 时，应用 IX 2.13(i) 的判据及 1.3 的明确计算。
$c>1$ 的情形用归纳法处理（参见 XVI 3.7）。设 $(Y,X)$ \Term{codimension}{余维} $c>1$。由定义，在 $X$ 上\Term{local-scheme}{局部}地可找到\Term{regular-triple}{正则三元组} $(Y,Z,X)$，使 $(Z,X)$ \Term{codimension}{余维} $1$，$(Y,Z)$ \Term{codimension}{余维} $c-1$。以
\[\begin{tikzcd}Y\arrow[r,"u"]\arrow[dr,"i"']&Z\arrow[d,"v"]\\&X\end{tikzcd}\]
记\Term{closed-immersion}{闭浸入}。有\Term{spectral-sequence}{谱序列}
\[E_2^{pq}=(R^pu^!)(R^qv^!)F\Longrightarrow(R^{p+q}i^!)F.\]
由归纳假设，$q\ne2$ 时 $(R^qv^!)F=0$，$q=2$ 时它\Term{locally-constant}{局部常值}等。再对 $u$ 应用归纳假设，得 $(p,q)\ne(2c-2,2)$ 时 $(R^pu^!)(R^qv^!)F=0$，而 $(p,q)=(2c-2,2)$ 时\Term{locally-constant}{局部常值}等，于是得到\Term{abutment}{终项} $(R^{p+q}i^!)F$ 的结果。
\end{proof}
\textbf{3.3。}现在还需把 3.2 中的\Term{locally-constant}{局部常值}\Term{sheaf}{层}确定到\Term{canonical}{典范}\Term{isomorphism}{同构}。对每个\Term{prescheme}{预概形} $S$，引入\Term{sheaf}{层} $(\mu_n^{\otimes r})_S$，即 $(\mu_n)_S$ 在常值环\Term{sheaf}{层} $\mathbf Z_S$ 上的 $r$ 重\Term{tensor-product}{张量积}。另令 $(\mu_n^{\otimes0})_S=(\mathbf Z/n)_S$。若 $n$ 在 $S$ 上\Term{invertible}{可逆}，它就是\Term{locally-constant}{局部常值} $n$ 阶循环群\Term{sheaf}{层}。有公式 $\mu_n^{\otimes r}\otimes\mu_n^{\otimes s}=\mu_n^{\otimes(r+s)}$；若 $S'\to S$，则 $(\mu_n^{\otimes r})_{S'}$ \Term{canonical}{典范}\Term{isomorphism}{同构}于 $(\mu_n^{\otimes r})_S$ 在 $S'$ 上的\Term{inverse-image}{逆像}。不致混淆时常省略 $S$。
设 $(Y,X)$ 为\Term{codimension}{余维} $1$ 的\Term{regular-pair}{正则对}，且 $Y$ 由方程 $x=0$ 定义。有\Term{sheaf}{层}的\Term{exact-sequence}{正合列}
\[0\longrightarrow(\mathbf G_m)_X\longrightarrow j_*(\mathbf G_m)_U\longrightarrow\mathbf Z_Y\longrightarrow0.\]
应用 1.3，得\Term{canonical}{典范}\Term{isomorphism}{同构}
\[(R^1j_*)\mu_n\simeq j_*(\mathbf G_m)_U/(j_*(\mathbf G_m)_U)^n\simeq(\mathbf Z/n)_Y,\tag{3.3.1}\]
\SourcePage{404}
它显然与 $x$ 的选择无关。因此有\Term{canonical}{典范}\Term{isomorphism}{同构}
\[\varphi:\mathbf Z/n\xrightarrow{\sim}(R^2i^!)\mu_n\simeq i^*(R^1j_*)\mu_n.\tag{3.3.2}\]
\begin{SGAresult}{定理}{3.4}
\footnote{依赖\Term{resolution-singularities}{奇点消解}（见引言）。本定理也可视为 3.2 与第 XVIII 讲中利用“\Term{local-fundamental-class}{局部基本类}”建立的变异理论的合取；该理论研究 $Y$ 在 $X$ 中\Term{local-complete-intersection}{局部完全交}时的\Term{homomorphism}{同态} $\varphi(Y,X)$（不一定\Term{bijection}{双射}）。}
设 $C$ 为\Term{excellent}{优良}\Term{equal-characteristic}{等特征}\Term{scheme}{概形}的\Term{regular-pair}{正则对}族。存在唯一函数 $\varphi$，对每个\Term{codimension}{余维} $c$ 的\Term{regular-pair}{正则对} $(Y,X)$ 给出\Term{sheaf}{层}的一个称为\Term{canonical}{典范}的\Term{isomorphism}{同构}
\[\varphi(Y,X)=\varphi:\mathbf Z/n\xrightarrow{\sim}\mathcal H_Y^{2c}(X,\mu_n^{\otimes c})=(R^{2c}i^!)\mu_n^{\otimes c},\]
满足如下条件。
(i) 若 $c=1$，且 $Y$ 在 $X$ 中由 $x=0$ 定义，则 $\varphi$ 为\Term{morphism}{态射} $(3.3\prime)$。\footnote{此处引注原文为 $(3.3\prime)$；前面的相应\Term{canonical}{典范}\Term{morphism}{态射}编号为 (3.3.2)。}
(ii) $\varphi$ 与\Term{smooth}{光滑}对的\Term{morphism}{态射}相容：若 $(Y',X')\to(Y,X)$ 为\Term{morphism}{态射}，$g:Y'\to Y$ 为诱导\Term{morphism}{态射}，则下列 $Y'$ \Term{sheaf}{层}图交换：
\[\begin{tikzcd}\mathbf Z/n\arrow[r,"{\varphi(Y',X')}"]&(R^{2c}i'^!)\mu_n^{\otimes c}\\g^*(\mathbf Z/n)\arrow[u]\arrow[r,"{g^*\varphi(Y,X)}"']&g^*(R^{2c}i^!)\mu_n^{\otimes c}\arrow[u].\end{tikzcd}\]
特别地，右侧竖\Term{arrow}{箭头}为\Term{isomorphism}{同构}。
(iii) \Term{transitivity}{传递性}：设 $(Z,Y,X)$ 为\Term{regular-triple}{正则三元组}，即\Term{regular}{正则}\Term{prescheme}{预概形}等的\Term{closed-immersion}{闭浸入}\Term{commutative-diagram}{交换图}
\[\begin{tikzcd}Z\arrow[r,"v"]\arrow[dr,"w"']&Y\arrow[d,"u"]\\&X.\end{tikzcd}\]
以 $a,b,c$ 分别记 $u,v,w$ 的\Term{codimension}{余维}，故 $c=a+b$。有\Term{isomorphism}{同构}
\[\varphi(Y,X)\otimes\mu_n^{\otimes b}:\mu_n^{\otimes b}\longrightarrow(R^{2a}u^!)\mu_n^{\otimes a}\otimes\mu_n^{\otimes b}\simeq(R^{2a}u^!)\mu_n^{\otimes c}.\]
\SourcePage{405}
\Term{spectral-sequence}{谱序列} $(R^pv^!)(R^qu^!)F\Rightarrow(R^{p+q}w^!)F$ 给出\Term{isomorphism}{同构}
\[\xi:(R^{2b}v^!)(R^{2a}u^!)\mu_n^{\otimes c}\xrightarrow{\sim}(R^{2c}w^!)\mu_n^{\otimes c}.\]
从而得到\Term{isomorphism}{同构}图
\[\begin{tikzcd}[column sep=large]\mathbf Z/n\arrow[r,"{\varphi(Z,Y)}"]\arrow[d,"{\varphi(Z,X)}"']&(R^{2b}v^!)\mu_n^{\otimes b}\arrow[d,"\eta"]\\(R^{2c}w^!)\mu_n^{\otimes c}&(R^{2b}v^!)(R^{2a}u^!)\mu_n^{\otimes c}\arrow[l,"\xi"]\end{tikzcd}\tag{3.4.1}\]
其中 $\eta=(R^{2b}v^!)(\varphi(Y,X)\otimes\mu_n^{\otimes b})$；\Term{transitivity}{传递性}断言此图交换。
\end{SGAresult}
\begin{proof}
存在唯一性对\Term{codimension}{余维} $c$ 归纳。$c=1$ 时，利用 (3.3.1) \Term{local-scheme}{局部}定义 $\varphi$。由于它与参数 $x$ 的选择无关，\Term{gluing}{粘合}后对每个\Term{codimension}{余维} $1$ 的 $(Y,X)$ 得到\Term{isomorphism}{同构} $\varphi$。由定义和 1.3，这样构造的 $\varphi$ 满足 (ii)。$c=1$ 时尚不涉及\Term{transitivity}{传递性}。
设对\Term{codimension}{余维}小于 $c$ 已证明存在唯一性，且 $c>1$。由归纳假设立即看出，在 3.4(iii) 的情形中，(3.4.1) 中已定义的三条\Term{arrow}{箭头} $\varphi(Z,Y),\xi,\eta$ 与\Term{regular-triple}{正则三元组}\Term{morphism}{态射}相容。令
\[\Psi(Z,Y,X)=\xi\eta\varphi(Z,Y).\]
则 $\Psi$ 也与三元组\Term{morphism}{态射}相容。若 $(Z,X)$ \Term{codimension}{余维} $c>1$，则在 $X$ 上\Term{local-scheme}{局部}可取三元组 $(Z,Y,X)$，满足 $0<\operatorname{codim}(Y,X)<c$。所以只须证明 $\Psi(Z,Y,X)$ 不依赖 $Y$。这将证明\Term{codimension}{余维} $c$ 时的唯一性，而存在性由\Term{gluing}{粘合}得出。
若任取 $(Z,Y,X)$，且 $\operatorname{codim}(Y,X)>1$，则在 $X$ 上\Term{local-scheme}{局部}可取\Term{regular-triple}{正则三元组} $(Y,W,X)$，使 $(W,X)$ \Term{codimension}{余维} $1$。通过一个留给读者的\Term{transitivity}{传递性}图追逐，便归结为在 $(Y,X)$ \Term{codimension}{余维} $1$ 时验证与 $Y$ 的独立性。
设 $(Z,Y_0,X),(Z,Y_1,X)$ 是两个这样的三元组。逐\Term{fiber}{纤维}验证即可，故可设 $X$ \Term{strictly-local}{严格局部}，且 $Y_i$ 由方程 $f_i=0$ 定义，$i=0,1$。令 $\overline X=\operatorname{Spec}\mathcal O_X[t]$，$\overline Z=Z\times_X\overline X$，以 $\overline Y$ 记由
\[(f_1-f_0)t+f_0=0\]
定义的闭子\Term{scheme}{概形}。显然 $\overline Y$ 在 $\overline X/X$ 的\Term{section}{截面} $\varepsilon_0:\{t=0\}$、$\varepsilon_1:\{t=1\}$ 上方各点\Term{regular}{正则}。令 $C$ 为 $\overline Y$ 不\Term{regular}{正则}的点所成的闭集（EGA IV 7.8.3(iv)）。
\SourcePage{406}
分别用 $\overline X-C,\overline Y-C,\overline Z-C\cap\overline Z$ 代替这些\Term{object}{对象}。于是 $(\overline Z,\overline Y,\overline X)$ 为\Term{regular-triple}{正则三元组}，\Term{section}{截面} $\varepsilon_i$ 给出三元组\Term{morphism}{态射}
\[(Z,Y_i,X)\longrightarrow(\overline Z,\overline Y,\overline X)\quad(i=0,1).\]
设 $F=(R^{2c}w^!)\mu_n^{\otimes c}$ 为 $Z$ 上的\Term{sheaf}{层}。它\Term{locally-constant}{局部常值}，并经任一\Term{isomorphism}{同构} $\Psi(Z,Y_i,X)$ \Term{isomorphism}{同构}于 $\mathbf Z/n$。容易看出，$\overline Z$ 上的 $\overline F=(R^{2c}\overline w^!)\mu_n^{\otimes c}$ 是 $F$ 的\Term{inverse-image}{逆像}。因 $\overline Z\to Z$ 的\Term{fiber}{纤维}\Term{connected-nonempty}{连通非空}，且该\Term{morphism}{态射}\Term{smooth}{光滑}，故为开映射，诱导\Term{morphism}{态射}
\[H^0(Z,F)\longrightarrow H^0(\overline Z,\overline F)\]
\Term{bijection}{双射}。因此图
\[\begin{tikzcd}& &H^0(Z,F)\\H^0(Z,F)\arrow[r]&H^0(\overline Z,\overline F)\arrow[ur,"\varepsilon_0^*"]\arrow[dr,"\varepsilon_1^*"']\\&&H^0(Z,F)\end{tikzcd}\tag{*}\]
中每个\Term{morphism}{态射}\Term{bijection}{双射}，故 $\varepsilon_0^*=\varepsilon_1^*$。给出 $F$ 与 $\mathbf Z/n$ 的一个\Term{isomorphism}{同构}，\Term{equivalence}{等价}于给出一个作为\Term{generator}{生成元}的\Term{global-section}{整体截面}（$1$ 的像）。设 $\alpha,\alpha_i$ 分别为 $\Psi(\overline Z,\overline Y,\overline X)$、$\Psi(Z,Y_i,X)$ 给出的 $\overline F,F$ 的\Term{global-section}{整体截面}。因 $\Psi$ 与\Term{morphism}{态射}相容，图 (*) 给出
\[\alpha_0=\varepsilon_0^*(\alpha)=\varepsilon_1^*(\alpha)=\alpha_1,\]
故 $\Psi(Z,Y_0,X)=\Psi(Z,Y_1,X)$。
\end{proof}
\SGAsection{4}{正则态射的局部无上同调性}
\begin{SGAresult}{定理}{4.1}
（\Term{locally-acyclic}{局部无上同调}性）\footnote{依赖\Term{resolution-singularities}{奇点消解}（见引言）。}设 $g:X'\to X$ 是\Term{excellent}{优良}\Term{equal-characteristic}{等特征}\Term{scheme}{概形}的\Term{regular}{正则}\Term{morphism}{态射}（EGA IV 6.8.1）。则对与\Term{characteristic}{特征}互素的 $n$，$g$ 是\Term{universal}{普遍}\Term{locally-acyclic}{局部无上同调}的。
\end{SGAresult}
\begin{proof}
假设在\Term{finite-type}{有限型}\Term{base-change}{基变换}下\Term{stable}{稳定}，由 XV 1.13(ii) 立即得出只须证明 $g$ 对 $n$ \Term{locally-acyclic}{局部无上同调}。应用 XV 1.17 的判据：取 $X'$ 的\Term{geometric-point}{几何点} $x'$，及 $X$ 中对应的\Term{geometric-point}{几何点} $x$。由\Term{regular}{正则}\Term{morphism}{态射}的定义立即得出，$g$ 诱导的这些点处\Term{strict-localization}{严格局部化}之间的\Term{morphism}{态射}仍\Term{regular}{正则}。于是可设
\SourcePage{407}
$X=\operatorname{Spec}A$ 和 $X'=\operatorname{Spec}A'$ \Term{strictly-local}{严格局部}，$g$ 是\Term{local-scheme}{局部}\Term{morphism}{态射}；据 XV 1.17，只须证明 $g$ 的\Term{geometric-fibre}{几何纤维}对 $n$ \Term{acyclic}{无上同调}。用 $X$ 的一个整闭子\Term{scheme}{概形}及 $X'$ 的对应闭子\Term{scheme}{概形}替换，便归结为几何\Term{generic-fiber}{泛纤维} $X'_x$ 的情形。取 $x=\operatorname{Spec}K$，其中 $K$ 是 $X$ 的\Term{rational-function-field}{有理函数域}的\Term{separable-closure}{可分闭包}。于是
\[x=\varprojlim_\alpha U_\alpha,\]
其中 $U_\alpha$ 遍历\Term{connected}{连通}\Term{etale}{平展} $X$ \Term{scheme}{概形}的一个\Term{filtered-category}{滤过范畴}，并可假设它们\Term{regular}{正则}（EGA IV 7.8.3(v)）。令 $B_\alpha$ 为 $A$ 在 $R(U_\alpha)$ 中的\Term{normalization}{正规化}，$U'_\alpha=X'\times_XU_\alpha$，$B'_\alpha=A'\otimes_AB_\alpha$。于是 $U_\alpha$ 是 $\operatorname{Spec}B_\alpha$ 的\Term{regular}{正则}\Term{open-subset}{开集}，$B_\alpha\to B'_\alpha$ \Term{regular}{正则}，所以由 (2.3)，
\[H^q(U_\alpha,\mu_n)\xrightarrow{\sim}H^q(U'_\alpha,\mu_n)\]
对每个 $q$ 成立。\Term{fiber}{纤维} $X'_x$ 是 $U'_\alpha$ 的\Term{limit}{极限}，故
\[\begin{aligned}H^q(X'_x,\mu_n)&=\varinjlim_\alpha H^q(U'_\alpha,\mu_n)\\&=\varinjlim_\alpha H^q(U_\alpha,\mu_n)\\&=H^q(x,\mu_n)=0\quad(q>0).\end{aligned}\]
\end{proof}
\begin{SGAresult}{推论}{4.2}
（\Term{regular}{正则}\Term{morphism}{态射}的\Term{base-change}{基变换}）设
\[\begin{tikzcd}X'\arrow[r,"g'"]\arrow[d,"f'"']&X\arrow[d,"f"]\\Y'\arrow[r,"g"']&Y\end{tikzcd}\]
为\Term{cartesian-diagram}{笛卡尔图}，$g:Y'\to Y$ 是\Term{excellent}{优良}\Term{equal-characteristic}{等特征}\Term{prescheme}{预概形}的\Term{regular}{正则}\Term{morphism}{态射}，$f$ \Term{quasi-compact}{拟紧}且\Term{quasi-separated}{拟分离}。则对每个与\Term{characteristic}{特征}互素的\Term{abelian-torsion-sheaf}{阿贝尔挠层} $F$，\Term{base-change}{基变换}\Term{morphism}{态射}（XII 4.2）
\[g^*(R^qf_*)F\longrightarrow(R^qf'_*)g'^*F\]
对每个 $q\ge0$ \Term{bijection}{双射}。
\end{SGAresult}
这由 4.1 及 XVI 1.1 得出。由 4.1 及 XV 1.17 立即得到：
\begin{SGAresult}{推论}{4.3}
设 $A$ 是\Term{strictly-local}{严格局部}\Term{excellent}{优良}\Term{equal-characteristic}{等特征}环，$f:X\to\operatorname{Spec}A$ 是\Term{morphism}{态射}。设 $A\to A'$ 是\Term{regular}{正则}\Term{homomorphism}{同态}，其中 $A'$ 也\Term{strictly-local}{严格局部}、\Term{excellent}{优良}、\Term{equal-characteristic}{等特征}。用撇号记该\Term{homomorphism}{同态}诱导的\Term{base-change}{基变换}。对 $X$ 上每个与 $A$ 的\Term{characteristic}{特征}互素的\Term{torsion-sheaf}{挠层} $F$，有
\[H^q(X,F)\xrightarrow{\sim}H^q(X',F')\]
\SourcePage{408}
对每个 $q\ge0$ 成立。
\end{SGAresult}
此推论特别适用于 $A'=\widehat A$ 为 $A$ 的\Term{completion}{完备化}的情形。（在这一情形，还可把“$A$ \Term{strictly-local}{严格局部}”的假设换为“$A$ 为 \Term{henselian-local-ring}{Hensel 局部环}”，这立即可见。）
\SGAsection{5}{有限性定理}
\begin{SGAresult}{定理}{5.1}
（\Term{finiteness}{有限性}）\footnote{依赖\Term{resolution-singularities}{奇点消解}（见引言）。}设 $f:X\to S$ 是\Term{excellent}{优良}\Term{equal-characteristic}{等特征}\Term{scheme}{概形}的\Term{finite-type}{有限型}\Term{morphism}{态射}，$F$ 是 $X$ 上与\Term{characteristic}{特征}互素的\Term{abelian-torsion-sheaf}{阿贝尔挠层}，并且是\Term{constructible}{可构造}\Term{sheaf}{层}。则对每个 $q$，$(R^qf_*)F$ 也\Term{constructible}{可构造}。
\end{SGAresult}
这推广了 XVI 5.1 的定理。证明几乎一样：以\Term{purity}{纯性}结果 3.2 代替 XVI 3.7。
\begin{SGAresult}{推论}{5.2}
设 $A$ 是\Term{excellent}{优良}\Term{strictly-local}{严格局部}\Term{equal-characteristic}{等特征}环，$f:X\to\operatorname{Spec}A$ 是\Term{finite-type}{有限型}\Term{morphism}{态射}。设 $F$ 是 $X$ 上与 $A$ 的\Term{characteristic}{特征}互素的\Term{constructible}{可构造}\Term{abelian-torsion-sheaf}{阿贝尔挠层}。则 $H^q(X,F)$ 对每个 $q$ 都是\Term{finite-group}{有限群}。
\end{SGAresult}
\SGAsection{6}{仿射态射的上同调维数}
设 $f:X\to Y$ 是\Term{excellent}{优良}\Term{prescheme}{预概形}的\Term{finite-type}{有限型}\Term{morphism}{态射}，$F$ 是 $X$ 上的\Term{sheaf}{层}。回顾在 $Y$ \Term{local-scheme}{局部}时已定义整数 $\delta(F)$（XIV 2.2）。因此对任意 $Y$，可把 $\delta$ 看成函数 $\delta(F,f)$，令
\[\delta(F,f,y)=\delta(F'),\tag{6.0}\]
其中 $Y'=\operatorname{Spec}\mathcal O_{Y,y}$，$X'\to Y'$ 由 $f$ 经\Term{base-change}{基变换} $Y'\to Y$ 得到，$F'$ 是 $F$ 在 $X'$ 上的\Term{inverse-image}{逆像}\Term{sheaf}{层}。
有下述推广 XVI 3.1 的结果。
\begin{SGAresult}{定理}{6.1}
\footnote{依赖\Term{resolution-singularities}{奇点消解}（见引言）。}设 $f:X\to Y$ 是\Term{affine}{仿射}\Term{finite-type}{有限型}\Term{morphism}{态射}，$Y$ 为\Term{excellent}{优良}\Term{equal-characteristic}{等特征}\Term{scheme}{概形}，$F$ 是 $X$ 上的\Term{abelian-torsion-sheaf}{阿贝尔挠层}。则 $R^qf_*F$ 在每个满足 $\delta(F,f,y)<q$ 的点 $y\in Y$ 为零。
\end{SGAresult}
显然\Term{equivalence}{等价}于 6.1 的\Term{local-scheme}{局部}\Term{form}{形式}如下。
\begin{SGAresult}{定理}{6.1 bis}
\footnote{依赖\Term{resolution-singularities}{奇点消解}（见引言）。}设 $A$ 是\Term{strictly-local}{严格局部}\Term{excellent}{优良}\Term{equal-characteristic}{等特征}环，$f:X\to Y=\operatorname{Spec}A$ 是\Term{affine}{仿射}\Term{finite-type}{有限型}\Term{morphism}{态射}，$F$ 是 $X$ 上的\Term{torsion-sheaf}{挠层}。则
\[H^q(X,F)=0\quad(q>\delta(F)).\]
\end{SGAresult}
\SourcePage{409}
回顾 X 5.1：对 $p$ 等于\Term{characteristic}{特征}的 $p$ \Term{torsion-sheaf}{挠层} $F$，这类断言没有什么意义。因此可设 $F$ 与\Term{characteristic}{特征}互素。有下述推论。
\begin{SGAresult}{推论}{6.2}
设 $A$ 是 Hensel \Term{excellent}{优良}\Term{equal-characteristic}{等特征}环，\Term{residue-field}{剩余域}为 $k$，$U\subset\operatorname{Spec}A$ 为\Term{affine}{仿射}\Term{open-subset}{开集}。设 $\ell\in\mathbf P$，则
\[\operatorname{cd}_\ell U\le\operatorname{cd}_\ell k+\dim A.\]
\end{SGAresult}
事实上，令 $\widetilde A$ 为 $A$ 的\Term{strict-localization}{严格局部化}；它是群为 $G=G(\bar k/k)$ 的无限 Galois \Term{etale}{平展}\Term{covering}{覆盖}。设 $\widetilde U$ 为由 $U$ 诱导的 $\operatorname{Spec}\widetilde A$ 的\Term{affine}{仿射}\Term{open-subset}{开集}。应用 Hochschild--Serre \Term{spectral-sequence}{谱序列}（VIII 8.4）
\[H^p(G,H^q(\widetilde U,\cdot))\Longrightarrow H^{p+q}(U,\cdot),\]
便归结为 $A=\widetilde A$ 的情形，此时是 6.1 bis 的特例。
沿用上述记号，以 $R(A)$ 记 $A$ 的\Term{rational-function-ring}{有理函数环}。有 $\operatorname{Spec}R(A)=\varprojlim_\alpha U_\alpha$，其中 $U_\alpha$ 是 $\operatorname{Spec}A$ 的\Term{affine}{仿射}\Term{open-subset}{开集}。故由取\Term{limit}{极限}理论，得到：
\begin{SGAresult}{推论}{6.3}
设 $A$ 是 Hensel \Term{excellent}{优良}\Term{equal-characteristic}{等特征}环，\Term{residue-field}{剩余域}为 $k$。设 $R(A)$ 为 $\operatorname{Spec}A$ 的\Term{rational-function-field}{有理函数域}。则对于 $\ell\in\mathbf P$，
\[\operatorname{cd}_\ell R(A)\le\operatorname{cd}_\ell k+\dim A.\]
若 $\ell$ 在 $k$ 中\Term{invertible}{可逆}，且 $\ell\ne2$ 或 $k$ 不可序，则取等号。
\end{SGAresult}
后一情形中的等号由 X 2.4 得出。最后，应用 X 4.2、X 4.4 得到：
\begin{SGAresult}{推论}{6.4}
设 $X$ 为 Noether \Term{excellent}{优良}\Term{equal-characteristic}{等特征}\Term{scheme}{概形}，$R(X)$ 为其\Term{rational-function-ring}{有理函数环}。设 $\ell\in\mathbf P$ 在 $X$ 上\Term{invertible}{可逆}，且 $\ell\ne2$ 或 $X$ 的任何\Term{residue-field}{剩余域}均不可序。则
\[\operatorname{cd}_\ell X\le\operatorname{cd}_\ell R(X)+\dim X\le2\operatorname{cd}_\ell R(X).\]
\end{SGAresult}
\textbf{6.1 的证明。}先作与 XIV 3 类似的归约。对 $\delta$ 归纳。显然，$\delta\le d$ 时的 6.1 \Term{equivalence}{等价}于 $\delta\le d$ 时的 6.1 bis。$\delta=0$ 时显然。用 6.1 bis 的\Term{form}{形式}处理 $\delta\le1$。因为 $F$ 是其\Term{constructible}{可构造}子\Term{sheaf}{层}的\Term{limit}{极限}（IX 2.9(iii)），可设 $F$ \Term{constructible}{可构造}。于是 $F$ 的\Term{support}{支集}包含在闭集 $Z\subset X$ 中，且每个 $z\in Z$ 都有 $\delta(z)\le1$（XIV 2.2）。$Z$ 的每个\Term{reduced}{约化}\Term{irreducible-component}{不可约分支} $Z_i$ 或为 $Y$ 的\Term{closed-point}{闭点} $y$ 上\Term{dimension}{维数}至多 $1$ 的\Term{affine}{仿射}\Term{finite-type}{有限型}\Term{scheme}{概形}，或在 $Y$ 的一个一维\Term{reduced}{约化}闭子\Term{scheme}{概形} $C$ 上\Term{quasi-finite}{拟有限}。后一情形，$C=\operatorname{Spec}B$，其中 $B$ 为一维整\Term{strictly-local}{严格局部}环。此时 $Z_i$ 是\Term{dimension}{维数}至多 $1$ 的整\Term{strictly-local}{严格局部}环的谱，或是 $B$ 的\Term{field-of-fractions}{分式域}的一个\Term{finite-extension}{有限扩张}域的谱。因此
\SourcePage{410}
由 X 2.3、IX 5.7，$Z_i$ 的\Term{cohomological-dimension}{上同调维数}在所有情形均至多为 $1$。有\Term{exact-sequence}{正合列}
\[0\longrightarrow F\longrightarrow\prod_i F_i\longrightarrow\varepsilon\longrightarrow0,\]
其中 $F_i$ 是 $F$ 在 $Z_i$ 上的\Term{restriction}{限制}，由此得出 $Z$ 的结果。
现在假设定理对 $\delta<d$ 已证明，来证明 $\delta(F)\le d$ 的情形，其中 $d\ge2$。
\begin{SGAresult}{引理}{6.5}
只须处理 $X=\mathbf E_Y^1=\operatorname{Spec}\mathcal O_Y[t]$ 为一维\Term{affine}{仿射}\Term{space}{空间}且 $Y$ \Term{strictly-local}{严格局部}的情形。
\end{SGAresult}
证明与 XIV 4.2 相同。显然可设 $Y$ \Term{strictly-local}{严格局部}且 $\delta(F)\le d$，须证明 $q>d$ 时 $H^q(X,F)=0$。可将 $X$ \Term{embedding}{嵌入} $\mathbf E_Y^N$，故只须处理 $X=\mathbf E_Y^N$。对 $N$ 归纳，考虑图
\[\begin{tikzcd}\mathbf E_Y^N\arrow[r,"g"]\arrow[dr,"f"']&\mathbf E_Y^{N-1}\arrow[d,"h"]\\&Y,\end{tikzcd}\]
使 $\mathbf E_Y^N$ 为 $\mathbf E_Y^{N-1}$ 上的一维\Term{affine}{仿射}\Term{space}{空间}。应用归纳假设及 Leray \Term{spectral-sequence}{谱序列}
\[H^p(\mathbf E_Y^{N-1},R^qg_*F)\Longrightarrow H^{p+q}(\mathbf E_Y^N,F),\]
知只须证明 $\delta(R^qg_*F)\le d-q$。
\SourcePage{411}
即若 $z\in\mathbf E_Y^{N-1}$ 满足 $\delta(z)>d-q$，则 $R^qg_*F$ 在 $z$ 为零。但关于 $N$ 的归纳假设适用于 $g$，故 $\delta(F,g,z)<q$ 时，该\Term{sheaf}{层}在 $z$ 为零。所以只须证明不等式
\[\delta(F,g,z)+\delta(z)\le\delta(F).\]
我们把这一乐趣留给读者。
\begin{SGAresult}{引理}{6.6}
为证明 $\delta\le d$ 时的 6.1，只须证明：设 $B$ 是\Term{dimension}{维数} $d$ 的\Term{complete}{完备}\Term{normal}{正规}\Term{equal-characteristic}{等特征}\Term{local-ring}{局部环}，$U\subset\operatorname{Spec}B$ 为\Term{affine}{仿射}\Term{open-subset}{开集}。假设存在 $x\in\operatorname{rad}B$ 在 $U$ 上\Term{invertible}{可逆}，则
\[H^q(U,\mathbf Z/\ell)=0\quad(q>d).\]
\end{SGAresult}
此引理也是推论 6.2 的特例。
\begin{proof}[证明（参见 XIV 4.4）]
由 6.5，只须处理 $Y$ \Term{strictly-local}{严格局部}且 $X=\mathbf E_Y^1$ 的情形。考虑图
\[\begin{tikzcd}\mathbf E_Y^1\arrow[r,"i"]\arrow[dr,"f"']&\mathbf P_Y^1\arrow[d,"g"]\\&Y,\end{tikzcd}\]
其中 $\mathbf P_Y^1$ 为\Term{projective-space}{射影空间}，$\mathbf E_Y^1$ 由它去掉无穷远\Term{section}{截面} $Y_\infty$ 得到。设 $F$ 是 $\mathbf E^1$ 上的\Term{torsion-sheaf}{挠层}，$\delta(F)\le d$。有\Term{spectral-sequence}{谱序列}
\[H^p(\mathbf P^1,R^qi_*F)\Longrightarrow H^{p+q}(\mathbf E^1,F).\]
$q>0$ 时 $R^qi_*F$ 集中于\Term{isomorphism}{同构}于 $Y$ 的\Term{section}{截面} $Y_\infty$。因 $Y$ \Term{strictly-local}{严格局部}，$p,q>0$ 时 $H^p(\mathbf P^1,R^qi_*F)=0$。又由 XII 5.3，$p>2$ 时 $H^p(\mathbf P^1,i_*F)=0$。因要证明 $n>d$ 时 $H^n(\mathbf E^1,F)=0$，且 $d\ge2$，只须证明
\[H^0(\mathbf P^1,R^qi_*F)=H^0(Y_\infty,R^qi_*F)=0\quad(q>d).\]
此群\Term{isomorphism}{同构}于 $R^qi_*F$ 在\Term{closed-fiber}{闭纤维}中 $\mathbf P^1$ 的无穷远点 $y_\infty$ 处的\Term{fiber}{纤维}。
由 VII 3.3、IX 2.9(iii)，只须处理\Term{constructible}{可构造} $F$。此时由 $\delta(F)\le d$，存在 $\mathbf P^1$ 的闭子\Term{scheme}{概形} $X$，\Term{dimension}{维数}至多 $d$，包含 $F$ 的\Term{support}{支集}。令 $\operatorname{Spec}A=\widetilde X$ 为 $X$ 在 $y_\infty$ 处的\Term{strict-localization}{严格局部化}，$\widetilde U=\widetilde X-\widetilde X\times_{\mathbf P^1}Y_\infty$，$\widetilde F$ 为 $F$ 在 $\widetilde U$ 上诱导的\Term{sheaf}{层}。注意 $\widetilde U$ 由 $\widetilde X$ 对元素 $x\in\Gamma(X,\mathcal O_X)$ \Term{localization}{局部化}得到。有
\[(R^qi_*F)_{y_\infty}\simeq H^q(\widetilde U,\widetilde F).\]
因此只须对 $\widetilde U$ 上每个\Term{torsion-sheaf}{挠层} $\widetilde F$ 证明 $q>d$ 时 $H^q(\widetilde U,\widetilde F)=0$。
现在取 $\varphi=\delta$ 应用 IX 5.6，知只须证明对于每个在 $\widetilde U$ 上\Term{finite}{有限}且整的 $V$ 和\Term{invertible}{可逆}的 $\ell$，
\[H^q(V,\mathbf Z/\ell)=0\quad(q>\delta(V)).\]
令 $\pi:\overline V\to V$ 为\Term{normalization}{正规化}。有\Term{exact-sequence}{正合列}
\[0\longrightarrow(\mathbf Z/\ell)_V\longrightarrow\pi_*(\mathbf Z/\ell)_{\overline V}\longrightarrow C\longrightarrow0,\]
其中 $\delta(C)<\delta(V)$。由此对 $\delta$ 归纳，只须取\Term{normal}{正规} $V$。
令 $B$ 为 $A$ 在 $V$ 的\Term{rational-function-field}{有理函数域}中的\Term{normalization}{正规化}。则 $B$ 是\Term{strictly-local}{严格局部}\Term{excellent}{优良}\Term{equal-characteristic}{等特征}环，$V$ 是 $\operatorname{Spec}B$ 的\Term{affine}{仿射}\Term{open-subset}{开集}，由对 $x$ \Term{localization}{局部化}得到。设 $\widehat B$ 为 $B$ 的\Term{completion}{完备化}，
\SourcePage{412}
$\widehat V$ 为诱导\Term{open-subset}{开集}（由 $\operatorname{Spec}\widehat B$ 对 $x$ 的像\Term{localization}{局部化}得到）。由 4.3，
\[H^q(V,\mathbf Z/\ell)\simeq H^q(\widehat V,\mathbf Z/\ell).\]
故只须证明
\[H^q(\widehat V,\mathbf Z/\ell)=0\quad(q>\delta(V)=\dim B=\dim\widehat B).\]
由归纳假设，$\delta(V)<d$ 时定理成立，故 6.1 bis 适用。因此只须处理 $\dim\widehat B=d$；因 $\widehat B$ \Term{normal}{正规}，正是引理的情形，结果得证。
\end{proof}
\textbf{6.7.0。}先在 $B$ \Term{characteristic}{特征}为 $p>0$ 时，加上额外假设证明 6.6。设 $A$ 是\Term{equidimensional}{等维}\Term{complete}{完备}\Term{equal-characteristic}{等特征}环，$k\subset A$ 是子域，使 $A/\operatorname{rad}A$ 为 $k$ 的\Term{finite-extension}{有限扩张}。称 $k$ \Term{algebra}{代数} $A$ 为\Term{formally-separable}{形式可分}的（formellement séparable），若存在形如 $k[[x_1,\ldots,x_n]]$ 的子 $k$ \Term{algebra}{代数}，使\Term{morphism}{态射} $\operatorname{Spec}A\to\operatorname{Spec}k[[x]]$ \Term{finite}{有限}，并在 $\operatorname{Spec}k[[x]]$ 的\Term{generic-point}{泛点}处\Term{etale}{平展}，其中 $x=(x_1,\ldots,x_n)$。
\begin{SGAresult}{引理}{6.7}
设 $k\subset B$ 为子域，$B/\operatorname{rad}B$ 在 $k$ 上\Term{finite}{有限}。在上述关于 $d$ 的归纳假设下，若 6.6 的假设成立，并额外满足下列任一条件，则 6.6 的结论成立：
(i) $k$ 的\Term{characteristic}{特征}为零。
(ii) 存在 $x\in\operatorname{rad}B$ 在 $U$ 上\Term{invertible}{可逆}，且 $\overline B=B/(x)$ 为\Term{formally-separable}{形式可分} $k$ \Term{algebra}{代数}。
\end{SGAresult}
\begin{proof}
取 $x_1=x,x_2,\ldots,x_d\in\operatorname{rad}B$，使它们生成一个以 $\operatorname{rad}B$ 为根的\Term{primary}{准素}\Term{ideal}{理想}。则 $B$ 是\Term{finite}{有限} $A$ \Term{algebra}{代数}，其中 $A=k[[x_1,\ldots,x_d]]$；$\overline B=B/x_1B$ 为 $\overline A=k[[x_2,\ldots,x_d]]$ \Term{algebra}{代数}。在 (ii) 中 $\overline B$ 解析\Term{separable}{可分}，故可选择 $x_2,\ldots,x_d$ 使 $\overline B$ 为泛\Term{etale}{平展}的 $k[[x_2,\ldots,x_d]]$ \Term{algebra}{代数}。于是 $B$ 也为泛\Term{etale}{平展} $A$ \Term{algebra}{代数}，当然这在 (i) 中也成立。
取 $b\in B$ 生成 $k((x))$ 的域\Term{group-extension}{扩张} $R(A)$。令 $\varepsilon\in k[[x]]$ 为 $b$ 在 $k[[x]]$ 上的首一\Term{irreducible}{不可约}方程的\Term{discriminant}{判别式}。在 (ii) 中，可选 $b$ 使 $x_1$ 不整除 $\varepsilon$。另写 $\operatorname{Spec}B-U=V(x_1)\cup Y$，其中 $Y$ 不含 $V(x_1)$ 的任何\Term{irreducible-component}{不可约分支}。令 $C=\Gamma(Y,\mathcal O_Y)$，其中 $Y$ 取诱导\Term{reduced-structure}{约化结构}。必要时改变 $x_2,\ldots,x_d$，可设 $V(x_d)$ 不包含在 $Y$ 在 $\operatorname{Spec}A$ 中的像内，且在 (ii) 中还可设 $V(x_d)$ 不包含在 $V(\varepsilon)$ 中。
\begin{SGAresult}{引理}{6.8}
令 $A^\circ=k[[x_1,\ldots,x_{d-1}]]\{x_d\}$ 为 $k[[x_1,\ldots,x_{d-1}]][x_d]$ 在 $x_1=\cdots=x_d=0$ 处的\Term{strict-localization}{严格局部化}。在上述条件下，存在 $A^\circ$ \Term{algebra}{代数}
\SourcePage{413}
$B^\circ$、\Term{affine}{仿射}\Term{open-subset}{开集} $U^\circ\subset\operatorname{Spec}B^\circ$ 和\Term{isomorphism}{同构} $A\otimes_{A^\circ}B^\circ\simeq B$，使 $U^\circ$ 在 $\operatorname{Spec}B$ 中诱导的\Term{open-subset}{开集}为 $U$。
\end{SGAresult}
此引理的证明类似 1.5。只须找 $A^\circ$ \Term{algebra}{代数} $B^\circ,C^\circ$ 及\Term{isomorphism}{同构} $A\otimes_{A^\circ}B^\circ\simeq B$、$A\otimes_{A^\circ}C^\circ\simeq C$；因为此时由 [3] 1.4，存在唯一（\Term{surjectivity}{满射}）\Term{morphism}{态射} $B^\circ\to C^\circ$ 诱导 $B\to C$。于是取 $U^\circ=\operatorname{Spec}B^\circ-\{\operatorname{Spec}C^\circ\cup V(x_1)\}$，立即看出 $U$ 由 $U^\circ$ 诱导。$C^\circ$ 的存在性以及 (ii) 中 $B^\circ$ 的存在性与 1.5 一样证明。(i) 中 $B^\circ$ 的存在性由 [3] 5.1 得出。（也可直接利用沿 $\{x_1=0\}$ 的\Term{ramified}{分歧}的 Abyankhar 引理以及\Term{descent}{下降}证明。）
现在可完成 6.7 的证明。再次应用 4.3，得到
\[H^q(U,\mathbf Z/\ell)\simeq H^q(U^\circ,\mathbf Z/\ell),\]
故只须证明右侧在 $q>d$ 时为零。$\operatorname{Spec}A^\circ$ 是 $k[[x_1,\ldots,x_{d-1}]][x_d]$ 上\Term{affine}{仿射}\Term{etale}{平展}\Term{scheme}{概形}的\Term{limit}{极限}，故 $\operatorname{Spec}B^\circ$ 是 $\operatorname{Spec}k[[x_1,\ldots,x_{d-1}]]$ 上\Term{finite-type}{有限型}\Term{scheme}{概形}的\Term{limit}{极限}，且其每个\Term{fiber}{纤维}\Term{dimension}{维数}至多为 $1$。由于 $U^\circ\subset\operatorname{Spec}B^\circ[1/x_1]$，有
\[U^\circ=\varprojlim_\alpha U^\circ_\alpha,\]
其中 $U^\circ_\alpha$ 在 $V^\circ=\operatorname{Spec}k[[x_1,\ldots,x_{d-1}]][1/x_1]$ 上\Term{affine}{仿射}且为\Term{finite-type}{有限型}，且其每个\Term{fiber}{纤维}\Term{dimension}{维数}至多为 $1$。
令 $f^\circ:U^\circ\to V^\circ$ 为结构\Term{morphism}{态射}。由于 $f^\circ_\alpha$ 的\Term{fiber}{纤维}\Term{dimension}{维数}均至多为 $1$，对 $v\in V^\circ$ 有
\[\delta(\mathbf Z/\ell,f^\circ_\alpha,v)\le\dim\mathcal O_{V^\circ,v}+1\le d-1\]
（因为 $\dim\mathcal O_{V^\circ,v}\le d-2$）。故关于 $d$ 的归纳假设适用，应用 6.1 得
\[R^qf^\circ_{\alpha*}(\mathbf Z/\ell)=0\text{ 在 }v\text{ 处，若 }q>\dim\mathcal O_{V^\circ,v}+1,\]
即若 $q>d-\dim\overline{\{v\}}$。因此该\Term{sheaf}{层}在 $v$ 非零蕴含 $\dim\overline{\{v\}}\le d-q$。所以对包含映射 $i:V^\circ\to\operatorname{Spec}k[[x_1,\ldots,x_{d-1}]]$（见 (6.0)），有
\[\delta(R^qf^\circ_{\alpha*}(\mathbf Z/\ell))\le d-q.\]
以 6.1 bis 的\Term{form}{形式}对 $i$ 应用归纳假设，得到
\[H^p(V^\circ,R^qf^\circ_{\alpha*}(\mathbf Z/\ell))=0\quad(p+q>d).\]
Leray \Term{spectral-sequence}{谱序列}
\[E_2^{pq}=H^p(V^\circ,R^qf^\circ_{\alpha*}(\mathbf Z/\ell))\Longrightarrow H^{p+q}(U^\circ_\alpha,\mathbf Z/\ell)\]
\SourcePage{414}
蕴含 $n>d$ 时 $H^n(U^\circ_\alpha,\mathbf Z/\ell)=0$。由于 $U^\circ$ 是 $U^\circ_\alpha$ 的\Term{limit}{极限}，同样得到
\[H^n(U^\circ,\mathbf Z/\ell)=0\quad(n>d).\]
这完成了 6.7 的证明，从而证明了\Term{characteristic-zero}{特征零}时的定理 6.1。
\end{proof}
\SGAsection{7}{仿射态射——证明的结束}
回顾，为在\Term{characteristic}{特征} $p>0$ 时完成 6.1 的证明（仍假设\Term{resolution}{消解}成立），由 6.6 已归结为对 $d\ge2$ 证明下述断言。
\begin{SGAresult}{引理}{7.1}
设 $B$ 是\Term{dimension}{维数} $d$ 的\Term{complete}{完备}\Term{normal}{正规}\Term{equal-characteristic}{等特征}\Term{local-ring}{局部环}，\Term{residue-field}{剩余域}\Term{separably-closed}{可分闭}。设 $U\subset\operatorname{Spec}B$ 为\Term{affine}{仿射}\Term{open-subset}{开集}，并存在 $x\in\operatorname{rad}B$ 在 $U$ 上\Term{invertible}{可逆}。则
\[H^q(U,\mathbf Z/\ell)=0\quad(q>d).\]
\end{SGAresult}
此外，由归纳假设，可设 6.1 对 $\delta<d$ 成立，且 7.1 在特例 6.7(ii) 中已得到证明。
先指出，$B$ \Term{normal}{正规}这一条件并不重要。事实上，设 $B$ 任意，$\pi:\overline B\to B$ 为其\Term{normalization}{正规化}，$\overline U\subset\operatorname{Spec}\overline B$ 为 $U$ 的\Term{inverse-image}{逆像}。有\Term{exact-sequence}{正合列}
\[0\longrightarrow(\mathbf Z/\ell)_U\longrightarrow\pi_*(\mathbf Z/\ell)_{\overline U}\longrightarrow C\longrightarrow0,\]
其中 $C$ 的\Term{support}{支集}\Term{dimension}{维数}小于 $d$，故 $q>d-1$ 时 $H^q(U,C)=0$。因 $H^q(U,\pi_*\cdot)\simeq H^q(\overline U,\cdot)$（VIII 5.6），可知
\[\bigl[H^q(U,\mathbf Z/\ell)=0\ (q>d)\bigr]\quad\Longleftrightarrow\quad\bigl[H^q(\overline U,\mathbf Z/\ell)=0\ (q>d)\bigr].\tag{7.2}\]
\textbf{7.3.0。}简要回顾\Term{formally-separable}{形式可分}性的一些所需性质，不给证明。读者可参阅 EGA IV 18.11.10、18.11.11。设 $A$ 是\Term{complete}{完备} Noether \Term{local-scheme}{局部} $k$ \Term{algebra}{代数}，\Term{residue-field}{剩余域}为 $k$ 的\Term{finite-extension}{有限扩张}。令 $A'$ 为\Term{local-ring}{局部环} $A\otimes_kk'$ 的\Term{completion}{完备化}，其中 $k'=k^{1/p}$。容易看出，\Term{group-extension}{扩张} $A\to A'$ 是\Term{radiciel}{根性}的（可能无限）；这是 Nagata 的著名例子。因此计算\Term{etale}{平展}\Term{cohomology}{上同调}时，用 $A'$ 替代 $A$ 无妨（VIII 1）。
\Term{dimension}{维数} $d$ 的\Term{equidimensional}{等维}\Term{algebra}{代数} $A$ \Term{formally-separable}{形式可分}，即为 $k[[x_1,\ldots,x_d]]$ 的\Term{finite}{有限}泛\Term{etale}{平展}\Term{group-extension}{扩张}，这一条件可用模 $\widehat\Omega^1_{A/k}$（EGA IV 18.11.10）或 \Term{jacobian-criterion}{Jacobian 判据}
\SourcePage{415}
（EGA $0_{\mathrm{IV}}$ 22.6、IV 7.1）\footnote{本研讨班的合编者未能阐明这条引文的意义：它所指的情形似乎很不同；也未能阐明随后断言的论证。A. G.}表述。由此看出，若 $A$ 不解析\Term{separable}{可分}，则 $k'$ \Term{algebra}{代数} $A'$ 不\Term{reduced}{约化}。例如，证明可归结到 $A$ 由一个方程 $f=0$ 给出的情形，其中 $f\in k[[x_1,\ldots,x_{d+1}]]$，并应用 \Term{jacobian-criterion}{Jacobian 判据}及下述引理；其有趣的证明留给读者。
\begin{SGAresult}{引理}{7.3}
设 $f(x_1,\ldots,x_n)\in k[[x_1,\ldots,x_n]]$ 为\Term{formal-power-series}{形式幂级数}。假设对每个 $i=1,\ldots,n$，存在\Term{invertible-element}{可逆元} $u_i\in k[[x_1,\ldots,x_n]]$，使 $u_if$ 为关于 $x_i^p$ 及 $x_1,\ldots,x_{i-1},x_{i+1},\ldots,x_n$ 的幂级数。则存在\Term{invertible-element}{可逆元} $v\in k[[x_1,\ldots,x_n]]$，使
\[vf\in k[[x_1^p,\ldots,x_n^p]].\]
\end{SGAresult}
由此容易推出，将 $(A,k)$ 换为 $(\overline{A'},k')$（其中 $\overline{A'}$ 为 $A'$ 的\Term{normalization}{正规化}）这一过程应用\Term{finite}{有限}次后，可归结到 $A$ 确实\Term{formally-separable}{形式可分}的情形。同样可证明，对给定 $x\in\operatorname{rad}A$，向\Term{algebra}{代数} $A$ 应用这一过程后，环 $(A/(x))_{\mathrm{red}}$ 将成为\Term{formally-separable}{形式可分}\Term{algebra}{代数}。
应用 VIII 1.1 及 7.2，得到：
\begin{SGAresult}{推论}{7.4}
只须在下列附加条件下证明 7.1：$B$ 包含子域 $k$，使 $B/\operatorname{rad}B$ 在 $k$ 上\Term{finite}{有限}，且 $k$ \Term{algebra}{代数} $(B/(x))_{\mathrm{red}}$ \Term{formally-separable}{形式可分}。
\end{SGAresult}
这一情形下，7.1 的证明归结为 6.7(ii) 已处理的特例。
沿用 7.1、7.4 的记号，令
\[D(x)=\sum_i r_iX_i\]
为 $x$ 的\Term{divisor}{除子}。于是每个 $X_i$ 是 $\operatorname{Spec}B$ 中\Term{codimension}{余维} $1$ 的\Term{irreducible}{不可约}闭子\Term{scheme}{概形}，在 $k$ 上\Term{formally-separable}{形式可分}，即 $B/J(X_i)$ 为\Term{formally-separable}{形式可分} $k$ \Term{algebra}{代数}；这也\Term{equivalence}{等价}于 $(B/(x))_{\mathrm{red}}$ \Term{formally-separable}{形式可分}。取 $y\in\operatorname{rad}B$，使
\[D(y)=\sum_i(r_i-1)X_i+\Gamma,\]
其中 $\Gamma$ 不含任何 $X_i$。考虑\Term{ideal}{理想} $(x,y)$ 的\Term{blow-up}{爆破}
\[f:Z\longrightarrow\operatorname{Spec}B.\]
环 $B[y/x]$ 及 $B[x/y]$ 的谱构成 $Z$ 的\Term{affine}{仿射}\Term{open-cover}{开覆盖}。
\SourcePage{416}
令 $B_1$ 为 $B[x/y]$ 在由 $(\operatorname{rad}B,x/y)$ 生成的\Term{maximal-ideal}{极大理想}处的\Term{strict-localization}{严格局部化}。对 $\operatorname{Spec}B[y/x]\to\operatorname{Spec}B$ 的\Term{closed-fiber}{闭纤维}的每个\Term{closed-point}{闭点} $Q$，以 $A_Q$ 记 $B[y/x]$ 在 $Q$ 处的\Term{strict-localization}{严格局部化}，以 $U_Q$ 记 $\operatorname{Spec}A_Q$ 中 $U$ 的\Term{inverse-image}{逆像}\Term{affine}{仿射}\Term{open-subset}{开集}。
\begin{SGAresult}{引理}{7.5}
为对环 $B$ 证明 7.1，只须证明对于 $q>d\ge2$，有
\[\text{(i)}\quad H^q(\operatorname{Spec}B_1[y/x],\mathbf Z/\ell)=0,\]
以及
\[\text{(ii)}\quad H^q(U_Q,\mathbf Z/\ell)=0\quad\text{对每个 }Q.\]
\end{SGAresult}
\begin{proof}
有\Term{open-immersion}{开浸入} $U\hookrightarrow\operatorname{Spec}B[1/x]\hookrightarrow\operatorname{Spec}B[y/x]$，从而有图
\[\begin{tikzcd}U\arrow[r,"i"]\arrow[dr,"k"']&\operatorname{Spec}B[y/x]\arrow[d,"j"]\\&Z.\end{tikzcd}\]
若 $Q$ 为 $\operatorname{Spec}B[y/x]\to\operatorname{Spec}B$ 的\Term{closed-fiber}{闭纤维}中的\Term{closed-point}{闭点}，由假设，$q>d$ 时 $(R^qi_*\mathbf Z/\ell)_Q=0$。对于 $\operatorname{Spec}B[y/x]$ 的每个不为该\Term{closed-fiber}{闭纤维}中\Term{closed-point}{闭点}的点 $P$，有 $\delta(F,i,P)<d$，其中 $F=(\mathbf Z/\ell)_U$（见 (6.0)）。因此对 $i$ 应用归纳假设及 6.1，得到对于\Term{morphism}{态射} $\operatorname{Spec}B[y/x]\to\operatorname{Spec}B$，
\[\delta(R^qi_*\mathbf Z/\ell)\le d-q.\]
以 6.1 bis 的\Term{form}{形式}对后一个\Term{morphism}{态射}应用归纳假设，得
\[H^p(\operatorname{Spec}B[y/x],R^qi_*\mathbf Z/\ell)=0\quad(q>0,\ p>d-q).\]
由 Leray \Term{spectral-sequence}{谱序列}
\[E_2^{pq}=H^p(\operatorname{Spec}B[y/x],R^qi_*\mathbf Z/\ell)\Longrightarrow H^{p+q}(U,\mathbf Z/\ell),\]
证明 7.1 归结为证明 $p>d$ 时 $H^p(\operatorname{Spec}B[y/x],i_*\mathbf Z/\ell)=0$。有\Term{exact-sequence}{正合列}
\[0\longrightarrow\mathbf Z/\ell\longrightarrow i_*(\mathbf Z/\ell)_U\longrightarrow C\longrightarrow0,\tag{7.6}\]
其中 $\delta(C)<d$。故再次由归纳，只须证明
\[H^p(\operatorname{Spec}B[y/x],\mathbf Z/\ell)=0\quad(p>d).\tag{7.7}\]
\SourcePage{417}
\Term{open-immersion}{开浸入} $j:\operatorname{Spec}B[y/x]\to Z$ 由 $Z$ 去掉“无穷远”的\Term{closed-subset}{闭子集} $C$ 得到；在 $Z$ 的另一个\Term{open-subset}{开集} $\operatorname{Spec}B[x/y]$ 中，$C=V(x/y)$。$R^qj_*\mathbf Z/\ell$ 集中于 $C$，而 $C\to\operatorname{Spec}B$ 为\Term{closed-immersion}{闭浸入}。因此 $p,q>0$ 时 $H^p(Z,R^qj_*\mathbf Z/\ell)=0$。由于 $Z\to\operatorname{Spec}B$ 的\Term{closed-fiber}{闭纤维}\Term{dimension}{维数}为 $1$，且该\Term{morphism}{态射}\Term{proper}{固有}，由 XII 5.3 bis 知 $p>2$ 时 $H^p(Z,j_*\mathbf Z/\ell)=0$。\Term{spectral-sequence}{谱序列}
\[E_2^{pq}=H^p(Z,R^qj_*\mathbf Z/\ell)\Longrightarrow H^{p+q}(\operatorname{Spec}B[y/x],\mathbf Z/\ell)\]
因而给出 $q>2$ 时的\Term{isomorphism}{同构}
\[H^q(\operatorname{Spec}B[y/x],\mathbf Z/\ell)\simeq H^0(Z,R^qj_*\mathbf Z/\ell)\simeq(R^qj_*\mathbf Z/\ell)_P,\]
其中 $P$ 是 $Z$ 的“无穷远”点。这些\Term{fiber}{纤维}正是（VIII 5）
\[(R^qj_*\mathbf Z/\ell)_P\simeq H^q(\operatorname{Spec}B_1[y/x],j_*\mathbf Z/\ell).\]
由 (7.7)，只须证明这些群在 $q>d\ge2$ 时为零。有一个与 (7.6) 类似的\Term{exact-sequence}{正合列}，结合归纳假设，知只须 $q>d$ 时 $H^q(\operatorname{Spec}B_1[y/x],\mathbf Z/\ell)=0$，引理得证。
\end{proof}
\textbf{7.5(i) 的证明。}给 $\operatorname{Spec}B[x/y]$ 的闭集 $C=V(x/y)$ 赋予诱导\Term{reduced-structure}{约化结构}。立即看出，$Z\to\operatorname{Spec}B$ 诱导的 $C$ 到 $\operatorname{Spec}B$ 的\Term{closed-immersion}{闭浸入}，将 $C$ 等同于闭子\Term{scheme}{概形} $X=\bigcup_iX_i$。由于
\[D(x)=\sum_i r_iX_i,\qquad D(y)=\sum_i(r_i-1)X_i+\Gamma,\]
元素 $x/y$ 在 $X_i$ 的\Term{generic-point}{泛点}附近生成\Term{ideal}{理想} $J(X_i)$；因为 $B$ \Term{normal}{正规}，那里的\Term{local-ring}{局部环}是\Term{discrete-valuation-ring}{离散赋值环}。因此 $\operatorname{Spec}B[x/y]\to\operatorname{Spec}B$ 在这些点附近是\Term{isomorphism}{同构}，且 $x/y$ 在 $C$ 的每个\Term{irreducible-component}{不可约分支} $C_i$（对应于 $X_i$）上以阶数 $1$ 消失。
由于闭集 $C$ \Term{lifting}{提升}到 $\operatorname{Spec}B_1$，用 $B_1$ 代替 $B[x/y]$ 后这仍成立，故 $B_1/(x/y)$ 是\Term{complete}{完备}环，且是\Term{formally-separable}{形式可分} $k$ \Term{algebra}{代数}。
为证明 7.5(i)，先应用 4.3，以\Term{completion}{完备化} $\widehat B_1$ 代替 $B_1$，这不改变 $B_1/(x/y)$。再由 (7.2)，可用其\Term{normalization}{正规化} $\overline B_1$ 代替 $\widehat B_1$。因 $B_1$（从而 $\widehat B_1$）在 $C_i$ 的\Term{generic-point}{泛点}处已经\Term{regular}{正则}，\Term{morphism}{态射} $\operatorname{Spec}\overline B_1\to\operatorname{Spec}\widehat B_1$ 在这些点是\Term{isomorphism}{同构}，
\SourcePage{418}
故 $\overline B_1/(x/y)$ 仍为\Term{formally-separable}{形式可分} $k$ \Term{algebra}{代数}。这样便归结为特例 6.7(ii)，得到结果。
\textbf{7.5(ii) 的证明。}取 $\operatorname{Spec}B[y/x]\to\operatorname{Spec}B$ 的\Term{closed-fiber}{闭纤维}上的点 $Q$。\Term{strict-localization}{严格局部化}可写成
\[A_Q=\varinjlim A,\]
其中 $A$ 遍历 $B[y/x]$ 上\Term{finite-type}{有限型}\Term{etale}{平展}环的一个\Term{filtered-category}{滤过范畴}。令 $U_A$ 为 $\operatorname{Spec}A$ 中 $U$ 的\Term{inverse-image}{逆像}\Term{affine}{仿射}\Term{open-subset}{开集}。因为 $x$ 在 $U_A$ 上\Term{invertible}{可逆}，\Term{morphism}{态射} $U_A\to\operatorname{Spec}B$ \Term{etale}{平展}，而 $U_A$ \Term{normal}{正规}。因此，若 $C$ 表示 $B$ 在 $A$ 的\Term{rational-function-field}{有理函数域} $R(A)$ 中的\Term{normalization}{正规化}，则有\Term{open-immersion}{开浸入}
\[U_A\longrightarrow\operatorname{Spec}C.\]
由 VII 5.7，只须证明，对每个 $A$，$q>d$ 时 $H^q(U_A,\mathbf Z/\ell)$ 在 $H^q(U_Q,\mathbf Z/\ell)$ 中的像为零。因 $x$ 在 $U_Q$ 上\Term{invertible}{可逆}，若 $f\in B$ 在 $B[y/x]$ 中能被 $x$ 整除，且 $f/x$ 在 $Q$ 不为零，则 $U_Q\to\operatorname{Spec}A$ 经 $\operatorname{Spec}A[1/f]$ \Term{resolution}{分解}。于是只须对某个适当的这种 $f$ 证明
\[H^q(U_{A,f},\mathbf Z/\ell)=0\quad(q>d),\]
其中 $U_{A,f}=U_A-V(f)$ 为 $\operatorname{Spec}A[1/f]$ 的\Term{affine}{仿射}\Term{open-subset}{开集}，也是 $\operatorname{Spec}C$ 的\Term{affine}{仿射}\Term{open-subset}{开集}。
取 $c\in k$ 使 $y/x+c$ 在 $Q$ 不为零，以 $J$ 记 $B$ 中闭集 $\Gamma\cap X$ 的\Term{ideal}{理想}。若
\[f\equiv y+cx\pmod{J^N},\qquad N>0,\tag{7.8}\]
则 $f$ 满足上述条件。事实上，在 $\operatorname{Spec}B[y/x]$ 中显然有 $V(J)=V(x)$。故 (7.8) 确实蕴含 $f$ 可被 $x$ 整除，从而
\[f/x\equiv y/x+c\pmod J.\]
因 $J$ 在 $Q$ 为零，而 $y/x+c$ 在 $Q$ 不为零，这证明了断言。
闭集 $\Gamma\cap X$ 在 $\operatorname{Spec}B$ 中\Term{codimension}{余维} $2$。因此可取满足 (7.8) 的 $f$，记为 $z$，使 $V(z)\cap V(y)$ 也\Term{codimension}{余维} $2$。再取 $x_3,\ldots,x_d$，使 $y,z,x_3,\ldots,x_d$ 为 $B$ 的\Term{system-of-parameters}{参数系}。于是 $B$，从而 $C$，为\Term{finite}{有限} $k[[y,z,x_3,\ldots,x_d]]$ \Term{algebra}{代数}。
令 $C^\circ$ 为 $k[[y,z,x_3,\ldots,x_d]]$ 在 $R(C)$ 中 $k((y,z,x))$ 的\Term{separable-closure}{可分闭包}内的\Term{normalization}{正规化}。于是 $C$ 为 $C^\circ$ 的\Term{radiciel}{根性}\Term{group-extension}{扩张}，而 $C^\circ$ 为 $k((y,z,x))$ 的泛\Term{etale}{平展}\Term{group-extension}{扩张}，其中 $x=(x_3,\ldots,x_d)$。
\SourcePage{419}
选取线性组合
\[f=az+by\qquad(a,b\in k),\]
使\Term{morphism}{态射} $\operatorname{Spec}C^\circ\to\operatorname{Spec}k[[z,y,x]]$ 在 $V(f)$ 的\Term{generic-point}{泛点}上方\Term{etale}{平展}，且在 $\operatorname{Spec}B[y/x]$ 中 $f/x$ 在 $Q$ 不为零。只须证明 $q>d$ 时 $H^q(U_{A,f},\mathbf Z/\ell)=0$，因而由 VIII 1.1，只须证明
\[H^q(U^\circ_{A,f},\mathbf Z/\ell)=0\quad(q>d),\]
其中 $U^\circ_{A,f}$ 为 $\operatorname{Spec}C^\circ$ 中对应于 $\operatorname{Spec}C$ 的\Term{open-subset}{开集} $U_{A,f}$ 的\Term{open-subset}{开集}。它是\Term{affine}{仿射}\Term{open-subset}{开集}，这可由\Term{descent}{下降}看出（EGA IV 2.7.1(xiii)）。由构造，$C^\circ$ 的元素 $f$ 在 $U^\circ_{A,f}$ 上\Term{invertible}{可逆}，且 $C^\circ/(f)$ 为\Term{formally-separable}{形式可分} $k$ \Term{algebra}{代数}。所以处于 6.7(ii) 的情形，这完成了 7.5(ii)，从而完成 6.1 的证明。
\SGAsection{}{参考文献}
\begin{enumerate}
\item Abyankhar, S.《\Term{characteristic}{特征} $p\ne0$ 的\Term{base-field}{基域}上\Term{algebra}{代数}曲面的\Term{local-scheme}{局部}一致化》（Local uniformization on algebraic surfaces over ground fields of characteristic $p\ne0$），Annals of Math.，第 63 卷（1963）；以及《\Term{embedding}{嵌入}\Term{algebra}{代数}曲面的\Term{resolution-singularities}{奇点消解}》（Resolution of singularities of embedded algebraic surfaces），Academic Press（1966）。
\item Abyankhar, S.《算术曲面的\Term{resolution-singularities}{奇点消解}》（Resolution of singularities of arithmetical surfaces），Arithmetical Algebraic Geometry，New York（1966）。
\item Artin, M.《Hensel 环上\Term{scheme}{概形}的\Term{etale}{平展}\Term{covering}{覆盖}》（Etale coverings of schemes over hensel rings），Amer. Journal Math.，第 88 卷第 4 期，第 915--934 页（1966）。
\item Hironaka, H.《\Term{characteristic-zero}{特征零}域上\Term{algebraic-variety}{代数簇}的\Term{resolution-singularities}{奇点消解} I、II》（Resolution of singularities of an algebraic variety over a field of characteristic zero I, II），Annals of Math.，第 79 卷第 1、2 期（1964）。
\item Weierstrass, K.《关于多变量解析函数的一些定理》（Einige auf die analytischen Functionen mehrerer Veränderlichen sich beziehende Sätze），Mathematische Werke von K. Weierstrass，第 2 卷，Berlin（1895）。
\item Bourbaki, N.《\Term{commutative-algebra}{交换代数}》（Algèbre commutative），第 3 章。
\end{enumerate}
\SourcePage{420}
\textbf{原文件插页。}$\mathrm{sga4\_19}$，第 385--412 页；《\Term{equal-characteristic}{等特征}\Term{excellent}{优良}\Term{prescheme}{预概形}的\Term{cohomology}{上同调}》，M. Artin。原页文字：“请单独编译此文件，并将它插入此处。”（法、英双语的同一句排版提示。）

